\documentclass[journal,twocolumn,9.5pt]{IEEEtran}
%\documentclass{elsart}
%%%%%%%%%%%%%%%%%%%%%%%%%%%%%%%%%%%%%%%%%%%%%%%%%%%
\usepackage{ifpdf}
\usepackage{graphicx,amssymb,cite,subfigure}
%%%%ieee comment
\usepackage{amsmath}
%%%%end ieee comment
%%%%%%%%%%%%%%%%%%%%%%%%%%%%%%%%%%%%%%%%%%%%%%%%%%%
\newtheorem{theorem}{Theorem}
\newtheorem{example}{Example}
\newtheorem{lemma}{Lemma}
\newtheorem{observation}{Observation}
\newtheorem{definition}{Definition}
%%%%ieee comment
%\newtheorem{proof}{Proof}
%%%%end ieee comment
\newtheorem{corollary}{Corollary}
\newcommand{\defeq}{\stackrel{\textrm{def}}{=}}
\newcommand{\wand}{\,\wedge\,}
\newcommand{\comom}[1]{\bar{#1}}
\renewcommand{\Pr}{P}
%%%%%%%%%%%%%%%%%%%%%%%%%%%%%%%%%%%%%%%%%%%%%%%%%%%
\renewcommand\floatpagefraction{.2}
\makeatletter
\def\elsartstyle{%
    \def\normalsize{\@setfontsize\normalsize\@xiipt{14.5}}
    \def\small{\@setfontsize\small\@xipt{13.6}}
    \let\footnotesize=\small
    \def\large{\@setfontsize\large\@xivpt{18}}
    \def\Large{\@setfontsize\Large\@xviipt{22}}
    \skip\@mpfootins = 18\p@ \@plus 2\p@
    \normalsize
}
\@ifundefined{square}{}{\let\Box\square}
\makeatother
\def\file#1{\texttt{#1}}
\pagestyle{plain}
%%%%%%%%%%%%%%%%%%%%%%%%%%%%%%%%%%%%%%%%%%%%%%%%%%%
\begin{document}

%\begin{frontmatter}
\title{CoMoM: Efficient Class-Oriented Evaluation of Multiclass Performance Models}
\author{Giuliano Casale, \IEEEmembership{Member,~IEEE}\thanks{Giuliano Casale ({\em casale@cs.wm.edu}) is with the College of William \& Mary, Department of Computer Science, 23187-8795, Williamsburg VA, USA.}}
\maketitle
\begin{abstract}
We introduce the {C}lass-{o}riented {M}ethod {o}f {M}oments (CoMoM), a new exact algorithm to compute performance indexes in closed multiclass queueing networks. Closed models are important for performance evaluation of multi-tier applications, but when the number of service classes is large they become too expensive to solve with exact methods such as Mean Value Analysis (MVA). CoMoM addresses this limitation by a new recursion that scales efficiently with the number of classes.

Compared to the MVA algorithm, which recursively computes mean queue-lengths, CoMoM carries on in the recursion also information on higher-order moments of queue-lengths. We show that this additional information greatly reduces the number of recursive steps needed to solve the model and makes CoMoM the best-available algorithm for networks with several classes. %For example, we show a model of a real J2EE application where CoMoM is five orders of magnitude faster than MVA and has memory occupation that is eight orders of magnitude smaller.

We conclude the paper by generalizing CoMoM to the efficient computation of marginal queue-length probabilities, which finds application in the evaluation of state-dependent attributes such as energy consumption or quality-of-service metrics.
\end{abstract}

\begin{keywords}
Multiclass applications, queueing networks, performance modeling, exact analysis, state-dependent attributes
\end{keywords}

\section{Introduction}
Closed product-form queueing networks are popular quantitative models for performance evaluation and capacity planning of multiclass systems, such as multi-tier applications processing several types of transactions\cite{BasCMP75,ChuAT97,MenA98,KouB03,UrgPSST05}. These systems are best modeled as closed due to the presence in real Web servers, application servers, and database servers of limits on the maximum number of concurrent user sessions~\cite{UrgPSST05}. We here focus on the probabilistic evaluation of models with several service classes, a case that cannot be handled efficiently with algorithms such as the widely-used Mean Value Analysis (MVA) \cite{ReiL80}, but which is often found in models of real applications. The main contribution of this paper is a new solution technique, called the Class-oriented Method of Moments (CoMoM), which can solve efficiently queueing networks with several classes.

The probabilistic evaluation of multiclass models is hard because the closed-form expressions of the state probabilities found in \cite{BasCMP75} include a {normalizing constant} that is very expensive to evaluate \cite{ReiK75b,ChaS80,ConG86,HarC02,Cas06b}. In particular, models with several classes are often infeasible since both MVA and multiclass Convolution \cite{ReiK75b} have requirements that grow exponentially in the number of classes. The {\sc Recal} class recursion \cite{ConG86} is tailored to models with several classes, but the limited scalability of this technique with respect to the population size makes it applicable only to models with few tens of requests. Recent studies on the analytical inversion of the generating function of the normalizing constant have led to the RGF class recursion \cite{HarC02,HarL04}, which is always more efficient than {\sc Recal}, but still remains too expensive for evaluating models with hundreds or thousands of requests. %Approximation methods for computing the normalization constants have also been recently explored \cite{RosTW94,ChoLW95}, but no accuracy bounds on the approximate solution are available.

%, which we show in case studies provide a much more accurate characterization of mean indexes for selecting the best cost-performance tradeoffs

In this paper, we dramatically improve over the requirements of the above algorithms by a new exact approach called the Class-oriented Method of Moments (CoMoM). Compared to the MVA, which recursively computes mean queue-lengths, CoMoM is based on a recursive evaluation of higher-order moments of queue-lengths. At each recursive step, CoMoM considers a new model with increased population and solves a linear system of equations to update the value of the higher-order moments of queue-lengths. This approach minimizes the number of steps required to complete execution since they grow only linearly with the total population size. Finally, CoMoM returns the normalizing constant from the computed queue-length moments.

By comparison with existing methods, we find that \textit{CoMoM is the best-available algorithm for evaluating queueing networks with several classes}. For example, we show a model of a real J2EE application where CoMoM is several orders of magnitude faster and more memory-efficient than MVA. We also compare CoMoM with the (queue-oriented) Method of Moments (MoM) \cite{Cas06b,Cas07}, which solves the model by recursively computing binomial moments of queue-lengths \cite{Cas07}. The analysis reveals that the combinatorial characteristics of the data structures used by CoMoM and MoM are similar, but CoMoM scales much better than MoM on models with several classes, while MoM is preferable on models with several queues.

%We also discuss in details the relations between CoMoM and the method of moments (MoM) presented in \cite{Cas06b}, showing that CoMoM is more efficient than the latter whenever the number of classes is greater than the number of servers, a situation that is typical in real studies, where the complexity of instrumenting real system typically limits the measurement process to a small number of servers compared to the number of applications running on them.

We conclude the paper by generalizing CoMoM to the efficient computation of marginal queue-lengths, which are needed to evaluate state-dependent indexes such as energy consumption or quality-of-service metrics. This analysis has never been performed efficiently by analytical methods, but we find that a generalization of CoMoM, called Pro-CoMoM, can perform the computation by an efficient recursive scheme. %We illustrate the technique on a case study related to energy management in Web servers.
%Energy management capabilities are rapidly spreading to all systems, e.g., to Ethernet communications \cite{ChrGSN08} and to RAID storage systems \cite{WanZL08}, and therefore have to be accounted in the capacity planning of real applications.  A correct evaluation of energy consumption with DVS technologies requires to quantify the concurrency levels at the resources because this technology is activated when a server is under-utilized compared to its maximum available capacity.

This paper is organized as follows. After giving required definitions in Section~\ref{SEC:background}, we introduce the basic ideas behind CoMoM in Section~\ref{SEC:illustrating} using an illustrative example. A general definition of the algorithm is given in Section~\ref{SEC:CoMoM}, followed in Section~\ref{SEC:linalg} by a detailed analysis of its numerical and linear algebraic properties. Computational costs and case studies validating the efficiency of CoMoM are given in Section \ref{SEC:efficiencyanalysis}. In Section \ref{SEC:statedepcosts}, we derive the Pro-CoMoM algorithm for computing marginal queue-lengths probabilities. A comparison of CoMoM with MoM is given in Section \ref{SEC:duality}. Conclusive remarks are reported in Section \ref{SEC:conclusion}. Appendix A describes a hybrid algorithm for singular models. We point to \cite{TR} for pseudo-codes and additional material.

% \subsection{Related Work}
% \label{SEC:related}
% We summarize recent contributions on computational methods for normalizing constants. We point to \cite{ConG86,BolGMT98} for reviews of established solution techniques, e.g., the Convolution algorithm \cite{Buz73,ReiK75b} or {\sc Recal} \cite{ConG86}.
%
% Starting from the seminal work of Bertozzi and McKenna \cite{BerM93}, several recent works have focused on computing normalizing constants using multidimensional generating functions. Choudhury, Leung and Whitt propose in \cite{ChoLW95} a numerical technique for approximately inverting the generating function. Low computational costs are obtained by truncating Euler summations arising in the generating function inversion and from exploiting sparsity of routing matrices. Harrison, Coury and Lee develop in \cite{HarC02,HarL04} the RGF exact algorithm, which computes the normalizing constant with a new recursion in terms of models with fewer classes. RGF derives from an explicit solution of the Cauchy's integral formula applied to the generating function of the normalizing constant and is more efficient than the {\sc Recal} class recursion \cite{ConG86,HarL04}.
%
% Parallel to the development of generating function methods, a number of works have applied sampling techniques to approximate the normalizing constant. Ross, Tsang and Wang develop in \cite{RosTW94,RosW97} a Monte Carlo summation and integration method for normalizing constants of multiclass models. The method is able to approximate models having a large number of stations, classes and customers, relies on the exponential sampling of the normalizing constant summation or integral formulas \cite{McKM84}, and returns confidence intervals on the computed estimates. An improvement of the Monte Carlo sampling approach is obtained by Tuffin in \cite{Tuf97} using a mixture of Monte Carlo and quasi-Monte Carlo methods. The paper also develops a variance reduction technique that applies to models in normal usage condition \cite{McKM84}.
%
% The solution of models with large populations has been also recently investigated using exact techniques. Casale proposes in \cite{Cas06b,Cas07} the method of moments (MoM), a recursive scheme based on a matrix recurrence equation involving simultaneously the recursive equations of LBANC \cite{ChaS80} and {\sc Recal} \cite{ConG86}. The method can be interpreted as a recursive computation of higher-order binomial moments of queue-lengths \cite{Cas07}. The relations between MoM and CoMoM are discussed in Section \ref{SEC:duality}.
% \begin{table}[ht!]
% \caption{Comparison of CoMoM and MVA on models with different number of queues and service classes. The population is partitioned equally among the classes. }
% \centering
% \renewcommand{\arraystretch}{1.1}
% \begin{tabular}{ccc| cc}
%   \hline
%   \multicolumn{3}{c|}{model parameters} &   \multicolumn{2}{c}{computational cost} \\
%   queues & classes & population & MVA & CoMoM \\
%   \hline
%   2 & 2 & 1000 & $10^{80}$ & $10^{10}$ \\
%   2 & 10 & 1000 & $10^{80}$ & $10^{10}$ \\
%   10 & 2 & 1000 & $10^{80}$ & $10^{10}$ \\
%   10 & 10 & 1000 & $10^{80}$ & $10^{10}$ \\
%   \hline
% \end{tabular}
% \end{table}
\section{Background}
\label{SEC:background}
A summary of the notation defined in this section is given in Table~\ref{TAB:notation}. We focus on the class of product-form models introduced in \cite{BasCMP75} and we consider a closed multiclass queueing network with constant population $\vec N\equiv (N_1,\ldots,N_r,\ldots,N_R)$, $N=\sum_{r=1}^R N_r$, where $R$ is the number of service classes and $N_r$ is the number of jobs (henceforth also called requests) of class~$r$. The $N$ jobs visit $M$ constant-rate queues according to a state-independent routing scheme.
The mean service demand of a class-$r$ job at queue $k$ is denoted by $D_{k,r}$, $1\leq k\leq M$, $1\leq r\leq R$, which is computed as a product between a mean service time and a mean visit ratio \cite{LazZGS84}. $Z_r$ is the think time of class-$r$ jobs before re-entering the network after completion; jobs in think state wait in a delay station indexed by $k=0$. Scheduling at the queues can be first-come first-served, processor sharing or last-come first-served preemptive resume \cite{BasCMP75}.
%, see \cite{BasCMP75} for related assumptions.

The equilibrium probability of the network being in state $\vec n=(n_{0,1},\ldots,n_{k,r},\ldots,n_{M,R})$, $\sum_{k=0}^M n_{k,r}=N_r$, $1\leq r\leq R$, is \cite{BasCMP75}
\begin{equation}
\label{EQU:bcmp}
\Pr(\vec n|\vec N)=\frac{1}{G(\vec N)}\left(\prod_{r=1}^R \frac{Z_{r}^{n_{0,r}}}{n_{0,r}!} \right) \left(\prod_{k=1}^M n_k!\prod_{r=1}^R \frac{D_{k,r}^{n_{k,r}}}{n_{k,r}!}\right),
\end{equation}
where $n_{k,r}$ is the number of class-$r$ jobs in station $k$, $n_k=\sum_{r=1}^R n_{k,r}$, and $G(\vec N)$ is a normalizing constant. % imposing~$\sum_{\vec n} \Pr(\vec n|\vec N)=1$. %; we point, e.g., to \cite{Lam83b} for formulas for computing mean performance indexes from normalizing constants.
\emph{Henceforth, we focus on the problem of computing $G(\vec N)$ and related quantities such as marginal queue-lengths obtained by summation of (\ref{EQU:bcmp}).}

Among the recurrence equations for computing normalizing constants presented in the literature, we use the  {\em convolution expression} (CE) \cite{ReiK75b,ChaS80}
\begin{equation}
\label{EQU:ce}
\textstyle G(\vec 1_k,\vec N)=G(\vec N)+\sum_{r=1}^R D_{k,r}G(\vec 1_k,\vec N-\vec 1_r),
\end{equation}
for arbitrary $1\leq k\leq M$, where the notation $\vec{1}_t$ stands for a vector composed of all zeros except for a one in the $t$-th position. Here, $G(\vec 1_k,\vec N-\vec{1}_r)$ is the normalizing constant of a model with a class-$r$ job less and augmented with a {\em replica} of queue $k$. A replica is an additional queue having service demands $D_{k,r}$, $1\leq r\leq R$, identical to the demands of queue $k$ in the original model. For consistency we define $G(\vec N)=G(\vec 0,\vec N)$, where $\vec 0=(0,0,\ldots,0)$. Equation (\ref{EQU:ce}) can be seen as a recursive formula for computing mean queue-lengths because the constants $G(\vec 1_k,\vec N)$ and $G(\vec 1_k,\vec N-\vec 1_r)$ may be interpreted as un-normalized queue-length values, see \cite{Lam83} and (\ref{EQU:meanxq}) given later. %In what follows, we instead use (\ref{EQU:ce}) for the first time in the computation of conditional moments of queue-lengths. Henceforth, we also refer to the normalizing constant $G(\vec v,\vec N)$ to denote a queueing network augmented with a set of queue replicas specified by the queue replica vector $\vec v$, where $\Delta m_k$ indicates the number of additional stations with demands identical to queue $k$ of the original model. For instance, $G(\vec 1_k+\vec 1_j,\vec N)=G(\vec N)$ is the normalizing constant of a model augmented with two queues identical to queue $k$ and $j$, respectively; in the case $\vec v=\vec 0$, we set $G(\vec 0,\vec N)=G(\vec N)$.

We also use the {\em population constraint}~(PC)
\begin{equation}
\label{EQU:pc}
\textstyle N_rG(\vec N)=Z_rG(\vec N-\vec 1_r)+\sum_{k=1}^M D_{k,r}G(\vec 1_k,\vec N-\vec 1_r),
\end{equation}
for $1\leq r\leq R$. Equation (\ref{EQU:pc}) is equivalent to the {\sc Recal} recurrence equation\footnote{Equation (\ref{EQU:pc}) can also be regarded as an un-normalized version of Little's Law\cite{Cas06b} when used to compute the number of jobs $N_Z$ in the delay server. This gives the expression $N_Z=N-\sum_{k=1}^M Q_{k,r}(\vec N)=Z_rX_r(\vec N)$. By (\ref{EQU:meanxq}) it is easy to verify that the last equivalence reduces to (\ref{EQU:pc}).} \cite{ConG86,Cas06b}.

We conclude observing that the mean class-$r$ throughput $X_r(\vec N)$ and the mean class-$r$ queue-length $Q_{k,r}$ of queue $k$ are computed from normalizing constants respectively as \cite{ReiK75b}
\begin{equation}
\label{EQU:meanxq}
X_{r}(\vec N)=\frac{G(\vec N-\vec 1_r)}{G(\vec N)}, \quad Q_{k,r}(\vec N)=D_{k,r}\frac{G(\vec 1_k,\vec N-\vec 1_r)}{G(\vec N)}.
\end{equation}
Response times and utilizations are easily obtained from throughputs and mean queue-lengths using Little's Law \cite{LazZGS84}.
 %Equations (\ref{EQU:ce})-(\ref{EQU:pc}) are essential to the definition of the CoMoM algorithm.
We point to comprehensive material in \cite{LazZGS84} for additional background on multiclass product-form queueing network models.
\begin{table}[ht!]
\renewcommand{\arraystretch}{1.05}
\centering
\caption{Summary of queueing network notation}
\begin{tabular}{|c|l|}
\hline
%~\\
{\em Symbol} & {\em Description}\\
\hline
$k,j$ & Queue indexes (range $1\ldots M$) \\
$r,s$ & Class indexes  (range $1\ldots R$)\\
${\vec 1}_t$ & Vector of zeros with a one in the $t$-th position\\
$D_{k,r}$ & Mean service demand of a class-$r$ request at queue $k$\\
$G(\vec N)$ & Normalizing constant\\
$G(\vec 1_k,\vec N)$ & Norm. constant of model with a replica of queue $k$ more\\
$M$ & Number of queues\\
$\vec N$ & Population vector \\
$N_r$ & Number of jobs (also called requests) of class $r$\\
$N$ & Total number of jobs in the network ($N=\sum_r N_r$)\\
$P(\vec n|\vec N)$ & Equilibrium probability of state $\vec n=(n_0,n_1,\ldots,n_M)$\\
$R$ & Number of service classes \\
$Z_r$ & Think time of class-$r$ requests\\
\hline
\end{tabular}
\label{TAB:notation}
\end{table}

\section{CoMoM Illustrative Example}
\label{SEC:illustrating}
To build intuition on the CoMoM approach and compare it with the standard MVA recursion\footnote{In order to perform a comparison using the normalizing constant approach, we here implicitly refer to the LBANC algorithm, i.e., the {\em un-normalized} version of MVA which computes normalizing constants instead of mean indexes (see \cite{ChaS80,Lam83}). }, we first show how to solve by CoMoM a small network with $R=2$ classes and $M=2$ queues. We focus on the computation of the normalizing constant for the population $(1000,1000)$ and we initially discuss the critical intermediate recursive step for $\vec N=(1000,1)$, which is the first population that cannot be evaluated by basic single class algorithms.  A graphical representation of the recursive structure of MVA in a neighborhood of $\vec N=(1000,1)$ is given in Figure~\ref{FIG:motivating}(a). MVA reaches $(1000,1)$ in a bottom-up approach, evaluating step-by-step models with populations summing to $N=N_1+N_2=1,\ldots,999,1000,1001$. The recursive step of MVA for $N=1001$ is depicted in Figure~\ref{FIG:motivating}(a), see the caption for a description. Figure~\ref{FIG:motivating}(b) illustrates the approach of the CoMoM recursion. CoMoM computes $\vec N=(1000,1)$ directly from the solutions of the single class models with populations $(1000,0)$, $(999,0)$, and $(998,0)$, which avoids evaluating almost half of the populations, i.e., the states $(1,1),(2,1),\ldots,(997,1)$ are all {\em skipped} by CoMoM. On larger models, while the number of states evaluated by MVA grows \emph{combinatorially} as the set of all possible populations component-wise less than or equal to $\vec N$, the set of states spanned by CoMoM grows \emph{linearly} with the total population size $N$, because the algorithm can recursively reapply the same step depicted in Figure~\ref{FIG:motivating}(b) for all populations. This efficient schema is made possible by simultaneously evaluating the CEs and PCs (\ref{EQU:ce})-(\ref{EQU:pc}) in a linear system of equations. %A comparison with the radically different queue-addition technique proposed in \cite{Cas06b}, that also employs CEs and PCs, is given in Section \ref{SEC:duality}.

In order to understand how the CoMoM technique performs a recursive step, let us consider the state $\vec N=(1000,1)$ and assume that we want to compute the set of normalizing constants
$\lambda(1000,1)=\{G(\vec 1_1,1000,1),G(\vec 1_2,1000,1), G(1000,1)\},$ i.e., the normalizing constant $G(1000,1)$ of the model and, according to (\ref{EQU:meanxq}), the un-normalized mean queue-lengths of stations $1$ and $2$.
Let us also assume to know from the previous recursive steps the similarly-defined vectors $\lambda(1000,0)$, $\lambda(999,0)$, $\lambda(998,0)$, which refer to models with a single class of jobs. These vectors are easily computed by efficient single-class algorithms. The CoMoM recursive step consists in the following recursive computation of the vectors $\lambda$.

\begin{figure}[ht!]
 \begin{center}
 \subfigure[MVA]{\includegraphics[width=0.4\columnwidth]{fig1a.eps}}
 \subfigure[CoMoM]{\includegraphics[width=0.4\columnwidth]{fig1b.eps}}
 \caption{{\bf Illustrative Example}.\em~Recursive structure of MVA (Figure~(a)) and CoMoM (Figure~(b)) in a neighborhood of $\vec N=(1000,1)$. Dashed states indicate populations of models solved previously by the recursion and not involved in the current recursive step; solid arrows indicate relations (e.g., (\ref{EQU:ce})-(\ref{EQU:pc})) that are exploited by the algorithm in the current step of the recursion.}
 \label{FIG:motivating}
 \end{center}
\end{figure}

\emph{Observation 1:} Using the PC in equation (\ref{EQU:pc}) for $r=2$, we can immediately compute $G(\vec N)=G(1000,1)$ from the normalizing constants in $\lambda(1000,0)$ which are assumed known. Thus, the unknowns that we still need to determine $\lambda(1000,1)$ are
\begin{align}
\label{EQU:exampleunk0}
\{G(\vec 1_1,1000,1),G(\vec 1_2,1000,1)\}\subset& \lambda(1000,1).
\end{align}

\emph{Observation 2:} We compute the constants (\ref{EQU:exampleunk0}) by the CEs (\ref{EQU:ce}). In fact, applying the PCs to compute $G(\vec 1_1,1000,1)$ or $G(\vec 1_2,1000,1)$ would require the availability of constants with two queue replicas that are not included\footnote{It is interesting to note that the MoM recursive step in \cite{Cas06b} differs from the CoMoM step described here because the $\lambda$ vector is modified to include also constants with two or more queue replicas. This makes the MoM recursion significantly different from the CoMoM recursion. We point to \cite{Cas06b,Cas07} and Section \ref{SEC:duality} for additional details.} in the vectors $\lambda$. Instead, the CEs can be used to evaluate $\lambda(1000,1)$ by including in the analysis also the constants in the vector $\lambda(999,1)$ which appear in the summation of (\ref{EQU:ce}) for $r=1$. Note that also $\lambda(999,1)$ is not known in advance and thus we have to expand the set of unknowns as
\begin{align}
\label{EQU:exampleunk1}
\{G(\vec 1_1,1000,1),G(\vec 1_2,1000,1)\}\subset& \lambda(1000,1),\\
\label{EQU:exampleunk2}
\{G(\vec 1_1,999,1),G(\vec 1_2,999,1)\}\subset& \lambda(999,1),
\end{align}
where $G(999,1)$ is omitted being easily computed by the PC for $r=2$ from the known constants in $\lambda(999,0)$.

\emph{Observation 3:}  The unknowns (\ref{EQU:exampleunk1}) are related to the unknowns (\ref{EQU:exampleunk2}) by two CEs for $k=1,2$, and the other constants which appear in these CEs are all known. The unknowns (\ref{EQU:exampleunk2}) are related to each other by the PC for $r=1$. Therefore, by extending the set of unknowns also to the $\lambda(999,1)$ vector we have moved from the initial problem (\ref{EQU:exampleunk0}) with $p=2$ unrelated unknowns to a new problem with $p^\prime=4$ unknowns (\ref{EQU:exampleunk1})-(\ref{EQU:exampleunk2}), but which are related by $q^\prime=3$ equations (\ref{EQU:ce})-(\ref{EQU:pc}). This reduces the degrees of freedom $df$ of the analysis from $df=p=2$ of (\ref{EQU:exampleunk0}) to $df=p^\prime-q^\prime=4-3=1$ of (\ref{EQU:exampleunk1})-(\ref{EQU:exampleunk2}). In other words, {\em the extension of the analysis to the vector $\lambda(999,1)$ in (\ref{EQU:exampleunk2}) has reduced the degrees of freedom in determining the normalizing constants}. This is a counter-intuitive property of multiclass models, since we are stating that increasing the difficulty of the analysis by adding to the unknowns also the constants in $\lambda(999,1)$ has the opposite effect of increasing our information about the system.

\emph{Observation 4:} Using the last observation, we can extend the analysis to include as unknown also the vector $\lambda(998,1)$ in addition to  $\lambda(1000,1)$ and  $\lambda(999,1)$. Using the same arguments given
above, we conclude that the new problem has $p^{\prime\prime}=6$ unknowns related by $q^{\prime\prime}=6$ linear equations (\ref{EQU:ce})-(\ref{EQU:pc}) and the degrees of freedom become $df=p^{\prime\prime}-q^{\prime\prime}=0$, which allows the exact computation of all unknowns by inverting a linear system of $q^{\prime\prime}$ equations (\ref{EQU:ce})-(\ref{EQU:pc}). Note that the degrees of freedom are reduced under the \emph{implicit assumption} that the $q^{\prime\prime}$ equations are linearly independent, a fact that in general depends on the specific values of the service demands $D_{k,r}$. We discuss cases of linear dependence between CEs and PCs in Section \ref{SEC:degenerate}, while in the rest of the paper we focus on non-degenerate cases.

Summarizing, starting from the vectors $\lambda(1000,0)$, $\lambda(999,0)$, and $\lambda(998,0)$, CoMoM has computed $\lambda(1000,1)$, $\lambda(999,1)$, and $\lambda(998,1)$. These vectors provide the new initial conditions for recursively applying the same scheme to models with more than one job in class $2$. With CoMoM we can easily reach the final model solution for the population $(1000,1000)$ after $1000$ recursive steps on class $2$. Conversely, MVA in Figure~\ref{FIG:motivating} still requires $990000$ evaluations to reach $(1000,1000)$. Thus, CoMoM can solve problems that are much harder than the ones that can be solved efficiently by the MVA algorithm. A general definition of the CoMoM algorithm for models with arbitrary populations, number of queues and classes is given below.

\section{The Class-Oriented Method of Moments}
\label{SEC:CoMoM}
Following the illustrative example, we give a general definition of CoMoM. The theorems and corollaries of this section are organized as follows. We first define the {\em basis} of CoMoM, which is a set of normalizing constants that summarizes the results of the past recursive steps performed by the algorithm and from which we can computer performance indexes by (\ref{EQU:meanxq}). We show later in the section that the basis of CoMoM can be seen as a set of higher-order moments of queue-lengths. In Theorem \ref{THM:CoMoM}, we explain how a basis may be computed recursively. Finally, in Corollary~\ref{COR:minimal} we give evidence that our definition of basis is parsimonious, i.e., CoMoM uses the minimum possible information needed to perform its simple recursion.
\begin{definition}{\em
Consider a model with $M$ queues, $R$ classes, and a population of $\vec N$ jobs, $N_r\geq 1$, $1\leq r\leq R$.
The {\em basis} ${\Lambda}(\vec N)$ of CoMoM is defined as the set
\begin{equation}
\label{EQU:lambda}
{\Lambda}(\vec N)={\lambda}(\vec 0)\cup \lambda(\vec 1_1)\cup \ldots \cup \lambda(\vec 1_M),
\end{equation}
where
$
\lambda(\cdot)=\{G(\cdot,\vec N-\vec n)\,|\, \textstyle \sum_{r=1}^{R-1} n_r\leq M {\,\rm \wedge\,} n_R=0\}
$
includes the normalizing constants of all models with population vectors having up to $M$ jobs less than $\vec N$ and where these jobs are chosen among the first $R-1$ service classes\footnote{By definition, for populations $\vec N-\vec n$ with one or more negative components, the normalizing constant is equal to zero; normalizing constants of models with population $\vec 0$ are instead equal to one. Thus, in models where some populations have $N_r\leq M$, some constants of (\ref{EQU:lambda}) are immediately computed as one or zero.}. }
\end{definition}

The above definition of $\Lambda(\vec N)$ always makes available enough CEs and PCs to reduce the degrees of freedom $df$ of the analysis to zero for any value of $M$, $R$ and $\vec N$. This statement is proved in the next theorem, which is the main result of this paper.

\begin{theorem}{\em
\label{THM:CoMoM}
The basis (\ref{EQU:lambda}) satisfies the linear matrix recurrence
\begin{equation}
\label{EQU:comom0}
{\bf \comom{A}}(\vec N){\Lambda}(\vec N)=\comom{\bf B}(\vec N){\Lambda}(\vec N-1_R),
\end{equation}
where ${\bf \comom{A}}(\vec N)$ and $\comom{\bf B}(\vec N)$ are square matrices of order
$\textstyle card(\Lambda(\vec N))=(M+1){M+R-1\choose M},$
defined by all equations (\ref{EQU:ce})-(\ref{EQU:pc}) that relate only
constants in ${\Lambda}(\vec N)$ and ${\Lambda}(\vec N-1_R)$.}
\end{theorem}
\begin{proof}
{\rm
The normalizing constants in ${\lambda}(\vec 0) \subseteq {\Lambda}(\vec N)$ are immediately computed from ${\Lambda}(\vec N-1_R)$ using the
PC for $r=R$. The remaining normalizing constants $\lambda(\vec 1_1)\cup \ldots \cup \lambda(\vec 1_M)$ are
$M\sum_{m=0}^M {m+(R-1)-1\choose m}=M{M+R-1\choose M},$
where we have observed that the sets $\lambda(\cdot)$ are uniquely defined by combinations with repetition of $m\in\{0,1,2,\ldots,M\}$ jobs chosen among $R-1$ classes. Observing that each PC requires to add a queue replica and remove a job with respect to the model in the left-hand side of (\ref{EQU:pc}), we can define $R-1$ PCs for each constant in $\{G( \vec 0,\vec N- \vec n)\,|\, \textstyle \sum_{r=1}^{R-1} n_r\leq M-1 {\,\rm \wedge\,} n_R=0\}\subset \Lambda(\vec N),$ which are thus in number
$(R-1){M+R-2\choose M-1}.$ For each normalizing constant in this subset, we can also define $M$ CEs, one for each possible queue $k$ replica to be added to the network, bringing the total number of available equations to
$(M+R-1){M+R-2\choose M-1}=M{M+R-1\choose M},$
which is exactly the cardinality of $\lambda(\vec 1_1)\cup \ldots \cup \lambda(\vec 1_M)$. Therefore, the matrices ${\bf \comom{A}}(\vec N)$ and ${\bf \comom{B}}(\vec N)$ defined by all the PCs and CEs relating only the constants in ${\Lambda}(\vec N)$ and ${\Lambda}(\vec N-1_R)$ are  both square of order $(M+1){M+R-1\choose M},$
and this also includes the PCs used to compute ${\lambda}(\vec 0) \subseteq {\Lambda}(\vec N)$.
}
%\qed
\end{proof}

We now examine in the next corollary the minimality of the definition of the basis ${\Lambda}(\vec N)$.
\begin{corollary}
\label{COR:minimal}
{\em
The basis ${\Lambda}(\vec N)$ has the minimal size to let ${\bf \comom{A}}(\vec N)$ and ${\bf \comom{B}}(\vec N)$ be always square on models with arbitrary number of queues $M$, number of classes $R$, and populations $\vec N$.}
\end{corollary}
\begin{proof}
Suppose to reduce the size of ${\Lambda}(\vec N)$ by defining
$
\lambda(\cdot)=\{G(\cdot,\vec N-\vec n)\,|\, \textstyle \sum_{r=1}^{R-1} n_r\leq C {\,\rm \wedge\,} n_R=0\}
$
for some integer $0\leq C<M$, then the size of the basis would be $(M+1){C+R-1\choose C},$ and the number of available equations would drop to
$(M+R-1){C+R-2\choose C-1}.$
Equating the number of unknowns and equations to have both ${\bf \comom{A}}(\vec N)$ and ${\bf \comom{B}}(\vec N)$ at least square, we get $(M+R-1){C+R-2\choose C-1}=(M+1){C+R-1\choose C},$
which simplifies to
${(M+R-1)}C=M{(C+R-1)}$ that is linear in $C$ and thus has the unique solution $C=M$ that proves the minimality of~(\ref{EQU:lambda}).
\end{proof}

The Class-oriented Method of Moments (CoMoM) algorithm follows from Theorem \ref{THM:CoMoM} by noting that, if ${\bf \comom{A}}(\vec N)$ is an invertible matrix, then ${\Lambda}(\vec N)$ is computed recursively as\footnote{Equation (\ref{EQU:comomrecursion}) holds also when $N_R=0$ and thus the right-hand side recursion must be performed on a class $r<R$. The only difference in this case is that ${\bf \comom{A}}(\vec N)$ and ${\bf \comom{B}}(\vec N)$ are defined without the PCs of class $R$ and some normalizing constants in the basis have negative populations (and thus these constants are set by definition equal to zero).}
\begin{equation}
\label{EQU:comomrecursion}
{\Lambda}(\vec N)={\bf \comom{A}}^{-1}(\vec N)\comom{\bf B}(\vec N){\Lambda}(\vec N-1_R).
\end{equation}
Note that (\ref{EQU:comomrecursion}) can be implemented in several ways, typically avoiding inversion of ${\bf \comom{A}}(\vec N)$, see Section \ref{SEC:linalg}.B for a discussion. Singularity of ${\bf \comom{A}}(\vec N)$ is also possible depending on the specific values of the demands $D_{k,r}$, we discuss the implications for CoMoM in Section \ref{SEC:degenerate}.
%Pseudo-code for computing the ${\bf \comom{A}}(\vec N)$ and  ${\bf \comom{B}}(\vec N)$ matrices from model parameters in the general case where $N_{r+1}=\ldots=N_{R}=0$ is given in the Appendix;
%CoMoM pseudo-code is given in Figure~\ref{ALG:comom}; the computation procedure for constructing $\comom{{\bf A}}(\vec N)$ and $\comom{{\bf B}}(\vec N)$ is given in Figure~\ref{ALG:pseudo}.
CoMoM pseudo-code and a computational procedure for generating the matrices $\comom{{\bf A}}(\vec N)$ and $\comom{{\bf B}}(\vec N)$ are reported in the technical report \cite{TR}. We also remark that an effective initialization of the recursion (\ref{EQU:comomrecursion}) can be done by an hybrid algorithm. That is, focusing on the un-normalized version of MVA which makes use of normalizing constants instead of mean-values, i.e., the LBANC algorithm \cite{ChaS80,Lam83}, we note that the set of normalizing constants in $\Lambda(N_1,0,\ldots,0)$ is exactly the set of constants computed by LBANC when evaluating the single-class populations $N_1,N_1-1,\ldots,N_1-M$. Thus, $\Lambda(N_1,0,\ldots,0)$ can be efficiently initialized by LBANC.

%\begin{figure}[ht!]
%\small
%\begin{center}
%\framebox
%{\begin{minipage}[h]{\columnwidth}
%\begin{tabbing}
%9\=99\=99\=99\=99\=99\=99\=999999999999999999999999999999999999999999999999\=99999999999999999999999999999999999999\=\kill
%Set constants with negative populations in $\Lambda(\vec 0)$ to zero \\
%Set constants with population $\vec 0$ in $\Lambda(\vec 0)$ to one \\
%\> {\bf for } $r$ {\bf from} $1$ {\bf to} $R$\\
%\> \> {\bf for } $n_r$ {\bf from} $1$ {\bf to} $N_r$\\
%\> \> \> $\vec n=(N_1,N_2,\ldots,N_{r-1},n_r,\vec 0,\ldots,\vec 0)$\\
%\> \> \> Compute $\comom{\bf A}(\vec n)$ and $\comom{\bf B}(\vec n)$\\
%\> \> \> $\Lambda(\vec n)=\comom{\bf A}^{-1}(\vec n)\comom{\bf B}(\vec n)\Lambda(\vec n-\vec 1_r)$\\
%\> \> {\bf end for}\\
%\> {\bf end for}\\
%\> Return $G(\vec N)$ from $\Lambda(\vec N)$
%\end{tabbing}
%\end{minipage}}
%\end{center}
%\caption{{\bf Basic CoMoM implementation.}
%%Pseudo-code for the computation of $\comom{\bf A}(\vec n)$ and $\comom{\bf B}(\vec n)$ is reported in the final Appendix.
%Significantly optimized implementations are also possible as we discuss later in Section \ref{SEC:btf}. In particular, one may compute $\comom{\bf A}(\vec n)$ and $\comom{\bf B}(\vec n)$ only at the beginning of a new cycle of the external for loop on $r$; a similar simplification applies also to the computation of $\comom{\bf A}^{-1}(\vec n)$ when the BTF of ${\bf A}$ is used.}
%\label{ALG:comom}
%\end{figure}
%
%\begin{figure}[h]
%\footnotesize
%\begin{center}
%\framebox
%{\begin{minipage}[h]{\columnwidth}
%\begin{tabbing}
%9\=99\=99\=99\=99\=99\=99\=999999999999999999999999999999999999999999999999\=99999999999999999999999999999999999999\=\kill
%%\> {{\bf Algorithm 1} \em - Computation of $\comom{\bf A}(\vec N)$ and $\comom{\bf B}(\vec N)$} \\
%\>	\>	$\comom{\bf A}(\vec N):=\{a_{ij}\}$ square matrix of order ${M+R\choose R}(M+1)$\\
%\>	\>	$\comom{\bf B}(\vec N):=\{b_{ij}\}$ square matrix of order ${M+R\choose R}(M+1)$\\
%\>	\>	Initialize all $a_{ij}$ and $b_{ij}$ to zero\\
%\>	\>	Initialize row counter $i$ to zero\\
%\>	\> {\bf for all} $ \vec n:=\{\sum_{r=1}^R  n_r=M,  n_r\geq 0\}$\\
%\>	\>	\> $ n_R:=0$; \\
%\>  \>  \>  {\bf if} $\sum_{r=1}^{R-1}  n_r<M$\\
%\>	\>	\>	\>  {\bf for } $k$ {\bf from } $1$ {\bf to } $M$\\
%\>  \>	\>	\>	$i:=i+1$;\\
%\>  \>	\>	\>	$j:=$ row of $G(\vec 1_k,\vec N- \vec n)$ in $\Lambda(\vec N)$;\\
%\>  \>	\>	\>	$a_{ij}:=1$;\\
%\>  \>	\>	\>	$j:=$ row of $G(\vec 0,\vec N- \vec n)$ in $\Lambda(\vec N)$;\\
%\>  \>	\>	\>	$a_{ij}:=-1$;\\
%\>	\>	\>	\>	\>  {\bf for } $s$ {\bf from } $1$ {\bf to } $R-1$\\
%\>  \>  \>	\>	\>	$j:=$ row of $G(\vec 1_k,\vec N- \vec n-\vec 1_s)$ in $\Lambda(\vec N)$;\\
%\>  \>  \>	\>	\>	$a_{ij}:=-D_{k,s}$;\\
%\>  \>  \>	\>	\>	{\bf end for}\\
%\>  \>	\>	\>	{\bf end for}\\
%\>	\>	\>	\>  {\bf for } $s$ {\bf from } $1$ {\bf to } $R-1$\\
%\>  \>	\>	\>	$i:=i+1$;\\
%\>  \>	\>	\>	$j:=$ row of $G(\vec 0,\vec N- \vec n)$ in $\Lambda(\vec N)$;\\
%\>  \>	\>	\>	$a_{ij}:=N_s- n_s$;\\
%\>  \>	\>	\>	$j:=$ row of $G(\vec 0,\vec N- \vec n-\vec 1_s)$ in $\Lambda(\vec N)$;\\
%\>  \>	\>	\>	$a_{ij}:=-Z_s$;\\
%\>	\>	\>	\>  \>	{\bf for } $k$ {\bf from } $1$ {\bf to } $M$\\
%\>  \>  \>	\>	\>	\>	$j:=$ row of $G(\vec 1_k,\vec N- \vec n-\vec 1_s)$ in $\Lambda(\vec N)$;\\
%\>  \>  \>	\>	\>	\>	$a_{ij}:=-D_{k,s}$;\\
%\>  \>	\>	\>	\>	{\bf end for}\\
%\>  \>	\>	\>	{\bf end for}\\
%\>	\>	\>	{\bf end if}\\
%\>	\>	\>	$i:=i+1$;\\
%\>	\>	\>	$j:=$ row of $G(\vec 0,\vec N- \vec n)$ in $\Lambda(\vec N)$;\\
%\>	\>	\>	$a_{ij}:=N_R$;\\
%\>	\>	\>	$j:=$ row of $G(\vec 0,\vec N- \vec n-\vec 1_R)$ in $\Lambda(\vec N)$;\\
%\>	\>	\>	$b_{ij}:=Z_R$;\\
%\>	\>	\>  {\bf for } $k$ {\bf from } $1$ {\bf to } $M$\\
%\>  \>	\>	\>	$j:=$ row of $G(\vec 1_k,\vec N- \vec n)$ in $\Lambda(\vec N-\vec 1-R)$;\\
%\>  \>	\>	\>	$b_{ij}:=D_{k,R}$;\\
%\>	\>	\>	{\bf end for}\\
%\>	\>	{\bf end for all}
%\end{tabbing}
%\end{minipage}}
%\end{center}
%\caption{{\bf Computation of $\comom{{\bf A}}(\vec N)$ and $\comom{{\bf B}}(\vec N)$ in CoMoM.} The indexes $i$ and $j$ in $a_{ij}$ and $b_{ij}$ denote rows and columns, respectively.}
%\label{ALG:pseudo}
%\end{figure}

%As justified by our remarks later in Section \ref{SEC:efficiency}, the computational cost of CoMoM is of the order of\footnote{To avoid the burden of the resulting formulas, these expressions ignore second-order and higher-order polylogarithmic factors in $N$, i.e., factors of the type $\log \log N, \ldots, \log \log \cdots \log N$. Details are given in Section \ref{SEC:efficiency}.} $$\left[{M+R-1 \choose M}(M+1)\right]^3 N^2 \log N$$ time and $$\left[{M+R-1 \choose M}(M+1)\right]^2 N \log N$$ space, which are on most models several orders of magnitude more efficient than MVA and class recursions such as {\sc Recal} \cite{ConG86}, which have requirements growing as $O(N^R)$ and $O(N^M)$ with the population size, respectively.
%We also point out that, as observed in general in \cite{Cas06b} for computations involving linear systems of CEs and PCs, the solution of the CoMoM linear system (\ref{EQU:comom0}) requires implementation in exact linear algebra (e.g., GNU GMP, Maple, Mathematica, MATLAB Symbolic Toolbox) because of well-known floating point-range exceptions in normalizing constant evaluations \cite{Lam82} and potential round-off error accumulations in the recursive solution of linear systems of equations (\ref{EQU:ce})-(\ref{EQU:pc})\cite{Cas06b}. Implementation in exact linear algebra is a definitive solution to these problems at the only expense of log-linear computational overheads in time and space which are discussed in Section~\ref{SEC:efficiency}.
%
%In the next subsections we study properties of the CoMoM algorithm including its probabilistic interpretation (Section \ref{SEC:probinterpret}), the behavior on degenerate models where $\comom{\bf A}(\vec N)$ is singular (Section \ref{SEC:degenerate}), and a BTF for $\comom{\bf A}(\vec N)$ with a related fine decomposition which both significantly reduce CoMoM requirements (Sections \ref{SEC:btf} to \ref{SEC:efficiency}).

\subsection{Probabilistic Interpretation}
\label{SEC:probinterpret}
We interpret CoMoM as a recursion on a set of higher-order moments of queue-lengths. We show in Theorem \ref{THM:moments} that these higher-order moments are conditioned on a group of jobs residing at a tagged station, which is consistent with previous work on the characterization of product-form models~\cite{Zah84}.  %Instead, the limited information available to the MVA imposes a different recursive structure that grows exponentially with the total number of classes. %For instance, in the illustrative example in Section \ref{SEC:illustrating} the higher-order moments in $\lambda(999,1)$ and $\lambda(998,1)$ carry the same information that the MVA would obtain by branching also on class~$1$, which is avoided by CoMoM since these vectors can be instead obtained by (\ref{EQU:comomrecursion}). % That is, this higher-order information is invested %As first observed by Zahorjan in \cite{Zah84}, the mean queue-length $Q_k(\vec N-\vec 1_r)$ of queue $k$ in a model with a class-$r$ job less, a quantity that is fundamental to the MVA recursive step, equals the mean queue-length $Q_k(\vec N\,|\,n_{k,r}\geq 1)$ conditioned on a class-$r$ job residing at $k$. Thus, MVA may also be interpreted as a solution of the model by reconstructing the mean job distribution across the network during residence times of tagged jobs. We now show that CoMoM implicitly generalizes this approach to the computation of higher-order conditional queue-lengths as seen by {\em groups} of jobs during their common residence times. In general, because of the additional factorial terms introduced in the probabilities (\ref{EQU:bcmp}) by a set of jobs conditioned to reside in a station, the CoMoM evaluation is not limited to basic {\em mean} queue-lengths, but also includes higher-order moments.
%\begin{theorem}\label{THM:probinter}{\em
%The basis $\Lambda(\vec N)$ is composed by a set of normalizing constants which is equivalent to the following conditional higher-order moments of queue-lengths:
%\begin{equation}\label{EQU:probint1}
%\frac{G(\vec N-\vec c_k)}{G(\vec N)}={\textstyle \sum}_{\vec c_k\leq \vec n_k} {\textstyle \prod}_{r=1}^R \alpha(\vec n_k, \vec c_k) \Pr(\vec n_k|\vec N),
%\end{equation}
%\begin{equation}\label{EQU:probint2}
%\frac{G(\vec 1_k,\vec N-\vec c_k)}{G(\vec N)}={\textstyle \sum}_{\vec c_k\leq \vec n_k} {\textstyle \prod}_{r=1}^R \beta(\vec n_k, \vec c_k) \Pr(\vec n_k|\vec N).
%\end{equation}
%where $\alpha(\vec n_k, \vec c_k)={(n_{k,r})_{c_{k,r}}}/({D^{c_{k,r}}_{k,r}(n_k)_{c_k}})$, $\beta(\vec n_k, \vec c_k)={(n_{k,r})_{c_{k,r}}}/({D^{c_{k,r}}_{k,r}(n_k)_{c_k-1}})$, $(x)_c=x(x-1)\cdots(x-c+1)$ is the falling factorial and $\Pr(\vec n_k|\vec N)$ is the marginal probability of queue $k$ in state $\vec n_k$. That is, (\ref{EQU:probint1})-(\ref{EQU:probint2}) are a set of higher-order inverse multinomial moments of the queue-lengths conditioned on the jobs in queue $k$ being $\vec c_k=(c_{k,1},c_{k,2},\ldots,c_{k,R})$.}
%\end{theorem}
%\begin{proof}{\rm The proof of the two formulas is similar, thus we focus on the more difficult (\ref{EQU:probint2}).
%Using the convolution definition of normalizing constants \cite{ReiK75b}, we write ${G(1_k,\vec N-\vec c_k)}$ as
%\begin{equation}
%\textstyle\sum_{\vec n_k\in \mathcal{S}(\vec c\,)}{G(-\vec 1_k,\vec N-\vec c_k-\vec n_k)}(n_k+1)!\left(\textstyle\prod_{r=1}^R {D_{k,r}^{n_{k,r}}}/{n_{k,r}!}\right),
%\end{equation}
%where the summation is on the set $\mathcal{S}(\vec c\,)=\{\vec n_k \,|\, 0\leq n_{k,r}\leq N_r-c_{k,r}, 1\leq r\leq R\}$, $\vec n_k$ is the population of the subnetwork composed by queue $k$ and its replica, the $(n_k+1)!$ factor also accounts for the convolution of the replica of queue $k$ \cite{ReiK75b}, and $G(-\vec 1_k,\vec N)$ is the normalizing constant of a model without queue~$k$. Rewriting the right-hand side of the above equation as
%\begin{multline}
%\sum_{\vec n\in \mathcal{S}(\vec c\,)} \left(\prod_{r=1}^R \frac{(n_{k,r}+c_{k,r})_{c_{k,r}}}{D^{c_{k,r}}_{k,r}(n_k+c_k)_{c_k-1}}\right)\cdot\\ \cdot\left(\prod_{r=1}^R \frac{G(-\vec 1_k,\vec N-\vec n_k-\vec c_k)}{G(\vec N)}\frac{ (n_k+c_{k})!D_{k,r}^{n_{k,r}+c_{k,r}}}{(n_{k,r}+c_{k,r})!}\right),
%\nonumber
%\end{multline}
% the result follows immediately by noting that the last parenthesis equals the probability $\Pr(\vec n_k|\vec N)$.}
%%\qed
%\end{proof}
%{\bf DISCUSSION, put in Appendix?}
%The
%According to Theorem \ref{THM:probinter}, the basis $\Lambda(\vec N)$, which is composed only by constants of the type $G(\vec 0,\vec N-\vec c_k)$ and $G(\vec 1_k,\vec N-\vec c_k)$ implicitly defines higher-order moments of queue-length defined on the states $\vec n_k\geq \vec c_k$ where a group of $\vec c_k$ jobs is conditioned to reside at queue $k$. Since the constants in $\Lambda(\vec N)$ are related to models with total number of jobs in $[N-M,N]$, we conclude that the group of jobs for which the higher-order conditional queue-lengths are evaluated by CoMoM is composed by $M$ jobs or less. Thus, this includes as special case the single residing job evaluated by the MVA.
%\newpage
%Compared to the MVA algorithm which recursively computes mean queue-lengths, CoMoM can be interpreted probabilistically as a recursive computation of higher-order conditional moments of queue-lengths.
\begin{theorem}
\label{THM:moments}
{\em
The basis $\Lambda(\vec N)$ is equivalent to a set of conditional higher-order moments of queue-lengths:
\begin{equation*}\label{EQU:probint1}\prod_{r=1}^R D^{c_{k,r}}_{k,r}\frac{G(\vec 0,\vec N-\vec c_k)}{G(\vec N)}=\sum_{\vec n:\vec n_k\geq \vec c_k} \prod_{r=1}^R \frac{(n_{k,r})_{c_{k,r}}}{(n_k)_{c_k}}\Pr(\vec n;\vec N),\end{equation*}
\begin{equation*}\label{EQU:probint2}\prod_{r=1}^R D^{c_{k,r}}_{k,r}\frac{G(1_k,\vec N-\vec c_k)}{G(\vec N)}=\sum_{\vec n:\vec n_k\geq \vec c_k} \prod_{r=1}^R \frac{(n_{k,r})_{c_{k,r}}}{(n_k)_{c_k-1}}\Pr(\vec n;\vec N),\end{equation*}
where $(x)_{c}=x(x-1)\cdots (x-c+1)$, $(x)_c=0$ if $x< c$, and $\vec c_k$ is a group of jobs conditioned to reside in queue $k$.
}\end{theorem}
% \begin{proof}{\rm The proof of the two formulas is similar, therefore we focus on the more complex of the two, i.e., equation (\ref{EQU:probint2}).
% Using (\ref{EQU:bcmp}), we expand the ratio ${G(1_k,\vec N-\vec c_k)}/{G(\vec N)}$ as
% \begin{equation}
% \label{TMP:0}
% \sum_{\vec n\in \mathcal{S}(\vec c\,)} \frac{G(-\vec 1_k,\vec N-\vec n_k-\vec c_k)}{G(\vec N)}\left((n_k+1)!\prod_{r=1}^R \frac{D_{k,r}^{n_{j,r}}}{n_{k,r}!}\right),
% \end{equation}
% where the summation is on the set $\mathcal{S}(\vec c\,)=\{\vec n | \sum_{j=0}^{M+1} n_{j,r}=N_r-c_{k,r}, 1\leq r\leq R\}$, the $n_k+1$ factor accounts for the replica of queue $k$, and $G(-\vec 1_k,\vec N)$ is the normalizing constant of a model without queue~$k$. Rewriting the right-hand side of (\ref{TMP:0}) as
% \begin{multline}
% \sum_{\vec n\in \mathcal{S}(\vec c\,)} \prod_{r=1}^R \frac{(n_{k,r}+c_{k,r})_{c_{k,r}}}{D^{c_{k,r}}_{k,r}(n_k+c_k)_{c_k-1}}\cdot\\ \cdot\left(\frac{G(-\vec 1_k,\vec N-\vec n_k-\vec c_k)}{G(\vec N)}\frac{ (n_k+c_{k})!D_{k,r}^{n_{k,r}+c_{k,r}}}{(n_{k,r}+c_{k,r})!}\right),
% \nonumber
% \end{multline}
% we note that the terms within parenthesis can be seen as the probability of the state $\vec n_k+\vec c_k$ in a model with population $\vec N$ and thus
% \begin{equation}
% \frac{G(1_k,\vec N-\vec c_k)}{G(\vec N)}=\sum_{\vec n\in \mathcal{S}(\vec 0\,):\vec n_k\geq \vec c_k} \prod_{r=1}^R \frac{(n_{k,r})_{c_{k,r}}}{D^{c_{k,r}}_{k,r}(n_k)_{c_k-1}}\Pr(\vec n;\vec N),
% \end{equation}
% which completes the proof.}
% %\qed
% \end{proof}
%, compared to the MVA approach, the CoMoM algorithm is able to exploit information on higher-order moments of queue-lengths that is not used by the MVA algorithm. This explains that
Due to limited space, we point to \cite{TR} for a proof of Theorem~\ref{THM:moments}. The intuition behind this finding is that, differently from the MVA that only evaluates mean values, CoMoM carries on in the recursion also higher-order information on the performance behavior of each server. The importance of higher-order moments is that CoMoM can use this extra information to decrease the degrees of freedom of the analysis and define by (\ref{EQU:comomrecursion}) a recursion that is scalable with the number of classes. {\em This is achieved by using the extra information to prevent CoMoM from stepping into the recursive branches of the classes} $r=1,\ldots,R-1$, as we have shown in the illustrating example in Section \ref{SEC:illustrating}. %The result clarifies that the CoMoM improved computational requirements with respect to MVA or multiclass Convolution originate from the availability at each recursive step of higher-order information about the model and not only of mean values as in the MVA.

\section{Numerical and Linear Algebraic Properties}
\label{SEC:linalg}
In this section, we discuss extensively the solution of the linear system~(\ref{EQU:comom0}). We first describe in Section \ref{SEC:btf} a permutation to block triangular form  for the coefficient matrix $\comom{\bf A}(\vec N)$ which improves scalability of CoMoM. Numerical properties of the recursive evaluation of (\ref{EQU:comom0}) are discussed in Section \ref{SEC:solution} where we also give recommendations on the best linear system solvers for CoMoM. Finally, in Section \ref{SEC:degenerate} we discuss degenerate models where (\ref{EQU:comom0}) is singular and define strategies to address these cases.

\subsection{Block Triangular Form and Fine-Grain Decomposition}
\label{SEC:btf}
We define a block triangular form (BTF) for $\comom{\bf A}(\vec N)$ in (\ref{EQU:comom0}) by observing that, while recurring on class $R$, the only entries of $\comom{\bf A}(\vec N)$ that change from $\comom{\bf A}(\vec N-\vec 1_R)$ are the $N_R$ coefficients in the PCs of class $R$. These PCs are needed only to compute the  constants in $\lambda(\vec 0)$, which suggests the following result.
\begin{theorem}
\label{THM:btf}
{\em The matrix $\comom{\bf A}(\vec N)$ can be permuted to the BTF
\begin{equation}
\label{EQU:btf}
\comom{\bf A}(\vec N)=
\begin{bmatrix}
\comom{\bf A}_{1,1} & \comom{\bf A}_{1,2} \\
{\bf 0} & N_R{\bf I}_{p} \\
\end{bmatrix}
\end{equation}
where $\comom{\bf A}_{1,1}$ is square of order $q={M+R-1\choose M}M$, the blocks $\comom{\bf A}_{1,1}$ and $\comom{\bf A}_{1,2}$ are independent of $N_R$, and ${\bf I}_{p}$ is the identity matrix of order $p={M+R-1\choose M}$ associated to the columns of $\lambda(\vec 0) \subset \Lambda(\vec N)$.
}\end{theorem}
\begin{proof}
{\rm The proof of the statement follows by permutation of rows and columns of $\comom{\bf A}(\vec N)$ according to the previous observation that only the PCs used to compute the unknowns in $\lambda(\vec 0)$ depend on $N_R$. The formulas for the cardinalities $p$ and $q$ follow immediately by the intermediate results in the proof of Theorem~\ref{THM:CoMoM}.}
%\qed
\end{proof}
Denote by $\Lambda_p(\vec N)$ the partition of $\Lambda(\vec N)$ associated to the unknowns in $\lambda(\vec 0)$, by $\Lambda_q(\vec N)$ the partition for the remaining constants, and by $\comom{\bf B}_{p}$ and $\comom{\bf B}_{q}$ the related sub-matrices of $\comom{\bf B}(\vec N)$ which are independent of $N_R$ as well. We rewrite the recursive step of CoMoM as
\begin{align}
\label{EQU:btfCoMoM}
\Lambda_q(\vec N)=&\comom{\bf A}^{-1}_{1,1}(\comom{\bf B}_{q}\Lambda(\vec N-1_R)-\comom{\bf A}_{1,2}\Lambda_p(\vec N-1_R)),\\
\Lambda_p(\vec N)=&N_R^{-1}\comom{\bf B}_{p}\Lambda(\vec N-1_R).
\end{align}
The main advantage of this form is that $\comom{\bf A}_{1,1}$ is smaller than $\comom{\bf A}(\vec N)$ and, more importantly, all matrices are independent of $N_R$, thus they can be computed and possibly factorized only once at the beginning of the iteration on class $R$.
% \subsection{Fine-Grain Decomposition}

Another optimization can be obtained by fine decomposition of $\comom{\bf A}_{1,1}$. Figure~\ref{FIG:A11CoMoM}, reported at the end of the paper, depicts $\comom{\bf A}_{1,1}$ for models with different number of service classes. We can see that also  $\comom{\bf A}_{1,1}$ admits a permutation to a fine-grain BTF.

\begin{theorem}
\label{THM:cblock}
{\em
The matrix $\comom{\bf A}_{1,1}$ admits the BTF
\begin{equation}
\label{EQU:cblock} \comom{\bf A}_{1,1}=
\begin{bmatrix}
\comom{\bf C}_{0} & \comom{\bf C}_{01} & {\bf 0}      & \ldots & {\bf 0} & {\bf 0}\\
{\bf 0}      & \comom{\bf C}_{1} & \comom{\bf C}_{12} & \ldots & {\bf 0} & {\bf 0}\\
{\bf 0}      & {\bf 0}      & \comom{\bf C}_{2} & \ldots & {\bf 0} & {\bf 0}\\
\vdots & \vdots & \vdots & \ddots & \vdots & \vdots\\
{\bf 0}      & {\bf 0}      & {\bf 0}      & \ldots & \comom{\bf C}_{\comom{H}-1} & \comom{\bf C}_{\comom{H}-1,\comom{H}} \\
{\bf 0}      & {\bf 0}      & {\bf 0}      & \ldots & {\bf 0} & \comom{\bf C}_{\comom{H}} \\
\end{bmatrix},
\end{equation}
where $\comom{H}=\min(M,R-1)$ and the macro-block
$\comom{\bf C}_{h}\equiv {\rm diag}(\comom{\bf C}^{(1)}_{h},\comom{\bf C}^{(2)}_{h},\ldots,\comom{\bf C}^{(t)}_{h},\ldots,\comom{\bf C}^{(T_h)}_{h})$, $0 \leq h \leq \comom{H}$, has $T_h={R-1 \choose h}$ micro-blocks $\comom{\bf C}^{(t)}_{h}$ of order ${M\choose h}M$.
}\end{theorem}
\begin{proof}
The basis $\Lambda(\vec N)$ is uniquely defined by the set of vectors $\vec n$ with $\sum_{s=1}^{R-1} n_s=M$, $n_s\geq 0$, which we partition according to the number and position of non-zeros in $\vec n$. We associate to each partition a value $h$ equal to the number of non-zeros in $\vec n$ associated to its elements; it is easy to verify that the range of $h$ is $0\leq h\leq \comom{H}$, $\comom{H}=\min\{M,R-1\}$. The macro-blocks $\comom{\bf C}_{h}$ are obtained by selecting the columns of $\comom{\bf A}_{1,1}$ and related rows associated to the normalizing constants with population $\vec N-\vec n$, where $\vec n$ has $h$ non-zeros. The matrix $\comom{\bf C}_h$ includes all micro-blocks $\comom{\bf C}_{h}^{(t)}$, which in number are equal to the way of assigning the positions of the $h$ non-zeros to the first $R-1$ classes of $\vec n$, i.e.,
${R-1 \choose h}$.
It is easy to verify that the cardinality ${M\choose h}M$ counts the number of feasible CEs and PCs in each micro-block and satisfies Vandermonde convolution \cite{Coh78} which states
$
{M+R-1 \choose R}R=\sum_{h=0}^{\comom{H}} {R-1\choose h} {M \choose h}M.
$
where the left-hand side is exactly the order of $\comom{\bf A}_{1,1}$.
\end{proof}
\begin{table}[t!]
\centering
\caption{Scalability of the Block Triangular Form (\ref{EQU:cblock}) of CoMoM}
\begin{tabular}{|p{6pt}p{6pt}|rr|rrr|}
\hline
&  &  &  &  &  & ~\\[-8pt]
\multicolumn{2}{|c|}{\sc Model} & \multicolumn{2}{|c|}{\sc Original Form (\ref{EQU:comom0})} & \multicolumn{3}{|c|}{\sc Block Triangular Form (\ref{EQU:cblock})}\\
\hline
&  &  &  &  &  & ~\\[-6pt]
&  & $\comom{\bf A}(\vec N)$ & $\comom{\bf A}(\vec N)$ & \textit{number of} & $\comom{\bf C}^{max}_{\comom{H}}$ & $\comom{\bf C}^{max}_{\comom{H}}$ \\
$M$ & $R$ &  \textit{order} & \textit{non-zeros} & \textit{micro-blocks} & \textit{order} & \textit{non-zeros} \\
\hline
&  &  &  &  &  & ~\\[-8pt]
2 & 2 & 9 & 22 &7 & 4 & 8 \\
2 & 4 & 30 & 95 &19 & 4 & 8 \\
2 & 8 & 108 & 397 &67 & 4 & 8 \\
2 & 10 & 165 & 626 &103 & 4 & 8 \\
\hline
&  &  &  &  &  & ~\\[-8pt]
4 & 2 & 25 & 76 &11 & 16 & 40 \\
4 & 4 & 175 & 777 &47 & 24 & 84 \\
4 & 8 & 1650 & 9494 &433 & 24 & 84 \\
4 & 10 & 3575 & 21870 &975 & 24 & 84 \\
\hline
&  &  &  &  &  & ~\\[-8pt]
8 & 2 & 81 & 280 &19 & 64 & 176 \\
8 & 4 & 1485 & 8481 &181 & 448 & 2464 \\
8 & 8 & 57915 & 487311 &6571 & 560 & 3640 \\
8 & 10 & 218790 & 2033603 &24829 & 560 & 3640 \\
\hline
&  &  &  &  &  & ~\\[-8pt]
10 & 2 & 121 & 430 &23 & 100 & 280 \\
10 & 4 & 3146 & 19071 &304 & 1200 & 6960 \\
10 & 8 & 213928 & 1988987 &19586 & 2520 & 20160  \\
10 & 10 & 1016158 & 10575708 & 92900 & 2520 & 20160 \\
\hline
\end{tabular}
\label{TAB:Areduction}
\end{table}
The form (\ref{EQU:cblock}) implies new computational savings, since (\ref{EQU:btfCoMoM}) can be computed by solution of a sequence of small linear systems with coefficient matrices $\comom{\bf C}^{(t)}_{h}$ that are much smaller than $\comom{\bf A}_{1,1}$. The effectiveness of the BTF is illustrated in Table~\ref{TAB:Areduction}. The table~reports order and number of non-zeros in $\comom{\bf A}(\vec N)$ and in the micro-block $\comom{C}_h^{(t)}$ of (\ref{EQU:cblock}) with largest order denoted by  $\comom{\bf C}^{max}_{{\comom{H}}}$. The properties of $\comom{\bf C}^{max}_{\comom{H}}$ characterize the hardest linear system solved by CoMoM during the evaluation of (\ref{EQU:btfCoMoM}) with BTF (\ref{EQU:cblock}); we also report the number of micro-block-level linear systems evaluated in (\ref{EQU:btfCoMoM})-(\ref{EQU:cblock}). We remark that the great majority of micro-blocks $\comom{\bf C}^{(t)}_{h}$ has much smaller order than $\comom{\bf C}^{max}_{\comom{H}}$. Table~\ref{TAB:Areduction} indicates that the improvements of the BTF are dramatic. Prohibitive linear systems with hundreds of thousands of equations are break into a long sequence of much smaller, feasible, linear systems.
We also observe that the BTF is more effective when $R>M$, e.g., the BTF is better for $M=8$ and $R=10$ than for $M=10$ and $R=8$. This is a desirable feature of the BTF because the focus of CoMoM is on models with large number of classes $R$, see Section~\ref{SEC:duality} for additional comments.

\subsection{Numerical Solution of CoMoM Linear Systems}
\label{SEC:solution}
We discuss the numerical evaluation of~(\ref{EQU:comom0}) in CoMoM with BTFs (\ref{EQU:btfCoMoM})-(\ref{EQU:cblock}), focusing on the properties of \textit{correctness} and \textit{efficiency} of the recursive solution.

{\em Correctness of CoMoM}. Regarding correctness, we argue that the solution of (\ref{EQU:comom0}), similarly to the MoM algorithm in \cite{Cas06b}, should be performed using {exact} linear algebra. This consideration arises from the fact that the solution of a {\em sequence} of linear systems of CEs and PCs using standard ({inexact}) linear algebra is subject to two numerical issues: floating-point range exceptions of normalizing constants \cite{Lam82} and potential round-off error accumulations while recurring on the sequence of linear systems~\cite{Cas06b}.
While floating-point range exceptions may be addressed using extended numerical precision (e.g, quad-precision arithmetic) or dynamic scaling \cite{Lam82}, we have observed that, similarly to MoM~\cite{Cas06b}, the error accumulation in CoMoM becomes uncontrollable with inexact algebra when the populations exceed more than a few tens of requests. The instability is visible also if each individual linear system (\ref{EQU:comom0}) is numerically well-conditioned (e.g., condition number $\kappa_{cond} < 10^2$) and the error accumulation can only be slowed, but not avoided, if extended precision arithmetic is used. For instance, Table~\ref{TAB:roundoff} shows an error accumulation example for a small model with $M=2$ queues and $R=2$ classes, see Section \ref{SEC:degenerate} for ${\bf \comom{A}}(\vec N)$ and $\Lambda(\vec N)$ structure. Here we have set $D_{1,1}=10 s$, $D_{2,1}=5 s$, $D_{1,2}=5 s$,  $D_{2,2}=9 s$, $Z_1=Z_2=0 s$, $N_1=150$ jobs and we have studied the error propagation in CoMoM solution as the population of class~$2$ grows. We have tested with linear system solution methods based on Gaussian elimination or iterative algorithms, including CG, BiCG, GMRES, and QMR implemented from the templates defined in~\cite{Bar94}. Table~\ref{TAB:roundoff} shows error accumulation for Gaussian elimination (implemented with LU back-substitution) and for the iterative algorithm that performed in the most stable~fashion, i.e., QMR with a maximum of $10,000$ iterations and convergence tolerance $10^{-13}$. In all four cases the numerical solution of (\ref{EQU:comom0}) becomes unstable when the population grows too large\footnote{Note that in Table~\ref{TAB:roundoff} the errors converge to $\Delta=1.00$ because CoMoM returns a normalizing constant that is several order of magnitude smaller than the exact one.}. Extended precision arithmetic helps in slowing down the error accumulation, but its effectiveness depends on the values of the service demands $D_{k,r}$. As a further remark, even if not immediately visible from the most significant digits of $G(\vec N)$, error accumulation starts as soon as the normalizing constant value exceeds the arithmetic range. This strongly suggests to use exact algebra to avoid introducing round-off errors in $G(\vec N)$.

\begin{table}[t]
\centering
\caption{Round-off Error Propagation with Inexact Linear Algebra}
\begin{tabular}{|c|cc|cc|}
\hline\\[-8pt]
      & \multicolumn{4}{|c|}{\em Relative Error Function: $\Delta=|1-G^{comom}(\vec N)/G^{exact}(\vec N)|$}  \\
\hline
      & \multicolumn{2}{|c|}{Gaussian Elimination} & \multicolumn{2}{|c|}{Iterative Approximation}  \\
$N_2$ & LU 32-bit &  LU 128-bit & QMR 32-bit & QMR 128-bit \\
\hline\\[-8pt]
$1$& $1.33\cdot 10^{-15}$ & $2.83\cdot 10^{-16}$ & $2.17\cdot 10^{-14}$& $2.83\cdot 10^{-16}$ \\
$25$& $1.73\cdot 10^{-11}$ & $6.46\cdot 10^{-16}$ & $7.84\cdot 10^{-8}$ & $9.86\cdot 10^{-16}$ \\
$50$& $2.98\cdot 10^{-5}$ & $1.94\cdot 10^{-16}$ & $0.15$ & $4.75\cdot 10^{-10}$ \\
$75$& $1.02$ & $2.56\cdot 10^{-15}$ & $1.00$ & $6.40\cdot 10^{-4}$ \\
$100$& $1.00$ & $2.06\cdot 10^{-9}$ & $1.00$ & $0.99$ \\
$125$& $1.00$ & $2.80\cdot 10^{-4}$ & $1.00$ & $1.00$ \\
$150$& $1.00$ & $1.00$ & $1.00$ & $1.00$ \\
\hline
\end{tabular}
\label{TAB:roundoff}
\end{table}

Following the last observation, we recommend for CoMoM a {definitive} solution to all {correctness} problems based on solving (\ref{EQU:comom0}) with exact linear algebra, which clearly imposes a trade-off with computational efficiency. Numerical stabilization using inexact linear algebra is a difficult open problem which leaves room for increasing the efficiency of CoMoM, see Section \ref{SEC:conclusion}. The use of exact algebra is motivated in the context of queueing network models by a fundamental observation: the computational overhead introduced by exact linear algebra can be upper bounded theoretically and it is found in Section \ref{SEC:efficiencyanalysis} to be much less than the gains resulting from the increased recursion efficiency of CoMoM compared to existing methods. Therefore, accepting exact linear algebra (\ref{EQU:comom0}) as a method to stabilize CoMoM grants the correctness of the results while leaving a consistent margin of computational gain with respect to the MVA and other methods. An upper bound on the computational overheads of exact algebra is given in the next result obtained in \cite{Cas06b}.

\begin{theorem}[Exact algebra overhead, \cite{Cas06b}]
{\em The computational overhead of adopting exact linear algebra in normalizing constant computations is upper bounded by the cost of adopting an arithmetic with operands having up to $n=N\lceil \log(D_{max}(N+M+R)) \rceil$ digits, where $D_{max}=\max_{k,r} \{D_{k,r},Z_r\}$ and $\lceil \cdot \rceil$ is the ceiling operator. For instance, when multi-precision arithmetic is used, the computational overheads of arithmetic operations grow as $O(n \log n \log \log n)$ \cite{SchS71}.}
\end{theorem}
\begin{proof} (Sketch of the proof, see \cite{Cas06b}.) The proof follows by observing that the number of digits required by the stabilization is upper bounded by $n=\lceil \log G \rceil$, where $G$ is the largest normalizing constant used in computations. The value of $n$ can be bounded as $n\leq n_{max}$, where $n_{max}=\lceil \log G^{max} \rceil$ is the number of digits of the normalizing constant of a model where we replace all $D_{k,r}$'s by $D_{max}=\max_{k,r} \{D_{k,r},Z_r\}$. Noting that $G^{max}$ refers to a model with balanced demands, it is easy to determine a closed-form expression for $G^{max}$ which returns $n_{max}=N\lceil \log(D_{max}(N+M+R)) \rceil$.
\end{proof}

{\em Efficiency of CoMoM}. From now on, we assume that CoMoM recursive evaluation of (\ref{EQU:comom0}) is implemented using exact algebra.
Among existing exact solution techniques for linear systems of equations, we recommend to use sparse finite-field linear algebra (FFLA) for the evaluation of ~(\ref{EQU:comom0}), see \cite{Vil07} and references therein for background. Compared to other exact methods, FFLA has minimal memory requirements thanks to the use of modulo arithmetic\footnote{We recall that in modulo arithmetic a linear system is decomposed into a set of finite-field linear systems over different modulos. These can be solved sequentially or in parallel using standard floating-point arithmetic and with memory costs similar to inexact linear algebra. The computed solutions are later assembled into the final result by the Chinese remainder theorem.} and time requirements grow only quadratically in the linear system order.

The best algorithm for solving (\ref{EQU:comom0}) using FFLA depends on the linear system size and sparsity. To achieve maximum scalability, one may use the Wiedemann algorithm\cite{LaMO90}, which is the standard solution method for large-scale systems in the high-performance Linbox library ({\em http://www.lingalg.org}), a C++ library resulting from a large collaborative research effort for the implementation of state-of-the-art FFLA algorithms. The Wiedemann method is a probabilistic algorithm for sparse linear systems over finite fields which may be seen as a FFLA counterpart of iterative algorithms such as conjugate gradient and Krylov subspace methods\cite{LaMO90}. The Wiedemann algorithm preserves the sparsity of the matrix $\comom{\bf A}(\vec N)$, which is crucial for CoMoM efficiency on large-scale models. For a linear system of order $n$ with $\omega$ non-zeros, the computational costs of Wiedemann algorithm grow as $O(n^2)$ in time and $O(n+\omega)$ in space \cite{LaMO90}, resulting in a much more efficient linear system solution than Gaussian elimination which requires $O(n^3)$ time and $O(n^2)$ space.
It should however be noted that for models that are not too large, e.g., $M<8 \wedge R<8$, the BTF (\ref{EQU:comom0}) results in a sequence of linear systems that have small order (few tens of elements) and where the coefficient matrix is not very sparse, see Table \ref{TAB:Areduction}. In these small linear systems, the cubic costs of LU factorization are typically small, the fill-ins do not grow too quickly, and LU back-substitution has quadratic costs which are generally much smaller than those of the Wiedemann algorithm. This can be explained by noting that the Wiedemann algorithm and iterative methods usually require some tens of iterations\footnote{This is consistent also with the recommendations of the Linbox library tutorials, where it is observed that Gaussian elimination should be used as the method of choice for solving small-scale or medium-scale sparse FFLA systems, whereas the Wiedemann algorithm is best for large systems.} before finding the exact solution of (\ref{EQU:comom0}). Based on these observations, the Wiedemann algorithm remains the method of choice for CoMoM on large-scale models, while Gaussian elimination is the best choice in all other cases. We show later in Case Study 3 of Section \ref{SEC:efficiencyanalysis} an example of model where the Wiedemann algorithm is much more scalable than Gaussian elimination in the evaluation of (\ref{EQU:comom0}).

{\em Computational Complexity}. The complexity formulas of the Wiedemann algorithm allow to estimate the computational costs of CoMoM under asymptotic growths of $M$, $R$, and $N$. In particular, from the intermediate results in theorems \ref{THM:CoMoM} and \ref{THM:cblock} which describe the number of CEs, PCs, and micro-blocks in $\comom{\bf A}(\vec N)$, it is not difficult to show that for (\ref{EQU:comom0}) with BTF (\ref{EQU:cblock}) the time requirements grow approximately~as $$\textstyle N^2\log(D_{max}(N+M+R))\sum_{h=1}^{\comom{H}} {R-1 \choose h} {M \choose h}^2 M^2,$$ which is the cost of solving $N$ micro-block level linear systems using the Wiedemann algorithm multiplied by the maximum numerical overhead of exact algebra\footnote{We refer to the worst-case scenario where the prime number $q$ defining the finite field is of the order of the normalizing constant, thus the overheads of exact arithmetic are $O(n \log n \log \log n)$ \cite{SchS71} with $n=N\lceil \log(D_{max}(N+M+R)) \rceil$. To simplify expressions, we do not write explicitly the polylogarithmic factor $\log \log n$ which has a very small impact on the complexity.}. Similarly, it is not difficult to show that space requirements grow approximately as
\begin{multline*}
\textstyle (\max_{1\leq h\leq \comom{H}} {M \choose h} M)+ M N\log(D_{max}(N+M+R)){M+R-1\choose M}\\ \textstyle +(M(R+1)+R(M+1)){M+R-2\choose M-1},
\end{multline*}
where the first term is the maximum Wiedemann algorithm space overhead, the middle term is a bound on the size of the basis $\Lambda(\vec N)$, and the last term is the storage cost of $\comom{\bf A}(\vec N)$ and $\comom{\bf B}(\vec N)$. Since existing methods such as the MVA, RECAL, and RGF grow in time and space requirements as $O(N^R)$ or $O(N^M)$, respectively, it is simple to conclude that for large enough populations CoMoM will be always more efficient than these methods since its complexity with respect to the population is $O(N^2\log N)$ in time and $O(N\log N)$ in space. In practice, the value of $N$ after which CoMoM becomes preferable is of the order of tens if the number of classes is large (e.g., $R> 3$). We point to Section \ref{SEC:efficiencyanalysis} for experiments confirming the efficiency of CoMoM. We conclude by remarking that the expression for computational complexity of Gaussian elimination, implemented with LU factorization and back-substitution, are worse than the ones of the Wiedemann algorithm: the time complexity requires to add a cubic term $\sum_{h=1}^{\comom{H}} {R-1 \choose h} {M \choose h}^3 M^3$ for LU factorization; the storage complexity requires to replace the memory overhead $(\max_{1\leq h\leq \comom{H}} {M \choose h} M)$ by a quadratic term $\sum_{h=1}^{\comom{H}} {R-1 \choose h} {M \choose h}^2 M^2$ that accounts for the fill-in resulting from the LU factorization. Nevertheless, the hidden constants of the asymptotic complexity notation make in practice Gaussian elimination more efficient than the Wiedemann algorithm for small- and medium-scale models.

%\begin{observation}
%{\em The numerical evaluation of (\ref{EQU:btfCoMoM}) with the Wiedemann algorithm has requirements that grow as $\textstyle  O\left({M+R-1\choose M}^2M^2\right)$ in time and as $\textstyle O\left((2MR+R-1){M+R-2\choose M-1}\right)$ in space.}
%\end{observation}
%\begin{proof}
%\end{proof}
%Experiments showing the scalability of CoMoM implemented with Wiedemann algorithm are reported in Section~\ref{SEC:efficiencyanalysis}.

%Iterative algorithms are even worse.
%Grwoth Number of digits normalization constant.
%These results are typical of models with same population values regardless of the values of the service demands $D_{k,r}$.

\subsection{Analysis of Degenerate Models}
\label{SEC:degenerate}
CoMoM can solve problems that are much harder than the ones that can be solved efficiently by the MVA, but the method requires independence between CEs and PCs. Occasionally, CEs and PCs are linearly dependent and this implies that ${\bf \comom{A}}(\vec N)$ is singular and the linear system (\ref{EQU:comom0}) does not have a unique solution. In this subsection, we discuss singularity conditions of ${\bf \comom{A}}(\vec N)$ and outline technical remedies to solve the queueing network model in this case. To simplify exposition, we illustrate singularity conditions using a small case study. We also explain how to generalize the techniques to models with larger number of queues or classes. Figure~\ref{FIG:singflowchart} summarizes the analysis of this section.

{\em Singularity Case Study}. Let us consider a small model composed by $M=2$ queues and $R=2$ classes. While recurring on class $2$, the left hand side ${\bf \comom{A}}(\vec N)\Lambda(\vec N)$ of the linear system (\ref{EQU:comom0}) is
\[
\renewcommand{\arraystretch}{1.1}
\arraycolsep 0.08em
\footnotesize
\begin{pmatrix}
    1&  \quad  0&    0&    0& -D_{1,1}&    0&    0&   0&    -1\\
    0&  \quad  1&    0&    0&    0& -D_{2,1}&    0&   0&    -1\\
    0&  \quad  0& -D_{1,1}& -D_{2,1}&    0&    0&  -Z_1& N_1-1&    0\\
    0&  \quad  0&    0& -D_{2,1}&    0&    1&    0&   -1&    0\\
    0&  \quad  0& -D_{1,1}&    0&    1&    0&    0&   -1&    0\\
    0&  \quad  0&    0&    0& -D_{1,1}& -D_{2,1}&    0&  -Z_1&   N_1\\
    0&  \quad  0&    0&    0&    0&    0&   N_2&    0&    0\\
    0&  \quad  0&    0&    0&    0&    0&    0&   N_2&    0\\
    0&  \quad  0&    0&    0&    0&    0&    0&    0&   N_2\\
\end{pmatrix}
\begin{pmatrix}
\small
\renewcommand{\arraystretch}{1.1}
\arraycolsep 0.1em
G(1_1,\vec N)\\
G(1_2,\vec N)\\
G(1_1,\vec N-\vec 2_1)\\
G(1_2,\vec N-\vec 2_1)\\
G(1_1,\vec N-\vec 1_1)\\
G(1_2,\vec N-\vec 1_1)\\
G(\vec N-\vec 2_1)\\
G(\vec N-\vec 1_1)\\
G(\vec N)\\
\end{pmatrix}
\]
where $\vec 2_1=2\cdot \vec 1_1$, rows 1 and 5 are CEs (\ref{EQU:ce}) for $k=1$, rows 2 and 4 are CEs for $k=2$, rows 3 and 6 are PCs (\ref{EQU:pc}) for $r=1$, and rows 7--9 are PCs (\ref{EQU:pc}) for $r=2$. The coefficients of (\ref{EQU:ce})-(\ref{EQU:pc}) not included above appear in the right hand side of (\ref{EQU:comom0}).
In the above example, the coefficient matrix has determinant $\det({\bf \comom{A}}(\vec N))=D_{1,1}D_{2,1}(D_{2,1}-D_{1,1})N_2^3,$
and, since $N_2\geq 1$, singularity arises if and only if $D_{k,1}=0$, for any $k=1,2$, or if $D_{1,1}=D_{2,1}$. We consider the two cases separately and show how we can solve these degeneracies.

\begin{figure}[t]
 \begin{center}
 \includegraphics[width=0.7\columnwidth]{fig2.eps}
 \caption{{\bf Addressing Singularity in CoMoM}.\em~ Summary of techniques proposed in Section \ref{SEC:degenerate} to address degenerate queueing network models.}
 \label{FIG:singflowchart}
 \end{center}
\end{figure}

{\em 1) Singularity Condition: Sparsity of Service Demands}. Let us assume without loss of generality that it is $D_{1,1}=0$. Then, column $3$ and column $5$ in ${\bf \comom{A}}(\vec N)$ are zero which implies $\det({\bf \comom{A}}(\vec N))=0$. This singularity can be prevented by removing from all bases the unknowns of columns $3$ and $5$, i.e. $G(1_1,\vec N-\vec 1_1)$  and $G(1_1,\vec N-\vec 2_1)$. The other unknowns can still be computed because they are independent of these two constants. In the general case, singularity due to sparsity of demands in larger models is addressed similarly by reducing the basis size.

{\em 2.a) Singularity Condition: Identical Queues}. The case $D_{1,1}=D_{2,1}$ is an example of degenerate demands belonging to a same service class. In this specific case, we can further distinguish two possibilities: (a) $D_{1,2}=D_{2,2}$ and queue $1$ and $2$ are identical for all demands; (b) $D_{1,2}\neq D_{2,2}$ which is discussed in the next singularity condition.

The singularity in subcase (a) can be handled easily because identical queues can be treated as one thanks to the following generalization of the PC \cite{Cas06b}
 \begin{equation}
 \label{EQU:genpc}
 \textstyle N_rG(\vec N)=Z_rG(\vec N-\vec 1_r)+\sum_{k=1}^M m_k D_{k,r}G(\vec 1_k,\vec N-\vec 1_r),
 \end{equation}
 for $1\leq r\leq R$, where $m_k\geq 1$ is the number of queues identical to queue $k$ (including queue $k$ itself). We point to \cite{Cas06b} for a proof of this equation. For instance, if $D_{1,1}=D_{2,1}$, we consider the model as composed by $M=1$ \emph{distinct} queues and $R=2$ classes. Thus, ${\bf \comom{A}}(\vec N)$ reduces to the non-singular matrix
\[
\renewcommand{\arraystretch}{1.1}
\arraycolsep 0.08em
\footnotesize
\begin{pmatrix}
   1& -D_{1,1}&   0&  -1\\
   0& -m_1 D_{1,1}& -Z_1&  N_1\\
   0&   0&  N_2&   0\\
   0&   0&   0&  N_2\\
\end{pmatrix}
\begin{pmatrix}
\small
\renewcommand{\arraystretch}{1.1}
\arraycolsep 0.1em
G(1_1,\vec N)\\
G(1_1,\vec N-\vec 1_1)\\
G(\vec N-\vec 1_1)\\
G(\vec N)\\
\end{pmatrix},
\]
where $m_1=2$ because there are two identical queues, row 1 is the CE for $k=1$, row 2 is the generalized PC (\ref{EQU:genpc}) for $r=1$, and rows 3 and 4 are the generalized PCs for $r=2$. Note that the linear system order decreases because for $M=1$ the basis $\Lambda(\vec N)$ is smaller than for $M=2$. We conclude this case by remarking that this solution approach based on (\ref{EQU:genpc}) generalizes to models of larger size by replacing sub-networks of identical queues by a single station and setting in (\ref{EQU:genpc}) the multiplicities $m_k$ accordingly.

{\em 2.b) Singularity Condition: Degenerate Demands for a Single Service Class}. The singularity condition of subcase (b), in which $D_{1,2}\neq D_{2,2}$, is instead more difficult, because the two queues are not identical and (\ref{EQU:genpc}) does not apply. However, this case can be handled effectively by changing the order in which we evaluate the service classes. That is, instead of processing first the class-$1$ populations $(1,0),(2,0),\ldots, (N_1,0)$, we first solve the model on the class-$2$ populations $(0,1),(0,2),\ldots, (0,N_2)$, which can be done  by single-class algorithms without incurring into singularity problems. After this step, CoMoM processes the class-$1$ populations $(1,N_2),(2,N_2),\ldots, (N_1,N_2)$. Since we have inverted the order in which we process the two classes, the coefficient matrix ${\bf \comom{A}}(\vec N)$ is now expressed as a function of $D_{k,2}$, $Z_2$, and $N_2$ instead of the corresponding parameters of class $1$ which were responsible for singularity of $\comom{\bf A}(\vec N)$. That is, while recurring on class $1$ to reach $(N_1,N_2)$, the parameters $D_{k,1}$, $Z_1$, and $N_1$ are now always included only in ${\bf \comom{B}}(\vec N)$ and hence cannot affect the singularity of ${\bf \comom{A}}(\vec N)$. This is sufficient to eliminate the singularity condition of identical demands in the linear system. Summarizing, this result is possible by the fact that the demands of the last processed class are {\em never} used as coefficients of ${\bf \comom{A}}(\vec N)$ because, from the structure of (\ref{EQU:ce})-(\ref{EQU:pc}), they appear only in ${\bf \comom{B}}(\vec N)$. This means that on larger models we can address {\em all} cases where some demands of a specific class are responsible for singularity by processing the population of this class as last.

{\em 3) Singularity Condition: Multiple Degeneracies}. Besides the conditions discussed for the case $M=2$ and $R=2$, other degenerate values of service demands arise in models with several queues and classes. For example, it is possible to show that in a model with $M=3$ queues and $R=3$ classes the coefficient matrix ${\bf \comom{A}}(\vec N)$ can become singular due to several degenerate demands, e.g., for $D_{2,2}=D_{3,2}D_{2,1}/D_{3,1}$. Since the number of singularity conditions increases with the model complexity it is possible to find cases where there are simultaneous degenerate demands belonging to different classes. These cases may not be addressable by changing the order of processing of service classes; thus, we define in Appendix A a hybrid MVA/CoMoM algorithm that can be used for models with multiple degeneracies. This is a general solution strategy where we focus on the independent equations in ${\bf \comom{A}}(\vec N)$. We stress that it is not possible to evaluate the singular linear system by approximate solvers because of the error accumulation issue, see Section \ref{SEC:solution}.

% Due to the large size of ${\bf \comom{A}}(\vec N)$ in these cases, we give an example focusing on the largest micro-block ${\bf \comom{C}}_h^{(t)}$ of the BTF (\ref{EQU:cblock}). Let us consider a model with $M=3$ queues and $R=3$ classes, the largest micro-block is $\comom{C}_{\comom{H}}^{1}$ which has structure
% \[
% \renewcommand{\arraystretch}{1.1}
% \arraycolsep 0.08em
% \footnotesize
% \begin{pmatrix}
%  -D_{1,2}& -D_{2,2}& -D_{3,2}&    0&    0&    0&    0&    0&    0\\
%     0& -D_{2,1}&    0&    0& -D_{2,2}&    0&    0&    1&    0\\
%  -D_{1,1}& -D_{2,1}& -D_{3,1}&    0&    0&    0&    0&    0&    0\\
%     0&    0&    0& -D_{1,1}& -D_{2,1}& -D_{3,1}&    0&    0&    0\\
%     0&    0&    0& -D_{1,2}& -D_{2,2}& -D_{3,2}&    0&    0&    0\\
%     0&    0& -D_{3,1}&    0&    0& -D_{3,2}&    0&    0&    1\\
%  -D_{1,1}&    0&    0& -D_{1,2}&    0&    0&    1&    0&    0\\
%     0&    0&    0&    0&    0&    0& -D_{1,1}& -D_{2,1}& -D_{3,1}\\
%     0&    0&    0&    0&    0&    0& -D_{1,2}& -D_{2,2}& -D_{3,2}\\
% \end{pmatrix},
% \]
% where rows including a $1$ element are CEs, while the remaining rows are PCs. The determinant of this micro-block is
% \begin{multline*}
% \det(\comom{C}_{\comom{H}}^{1})=\Delta_a\Delta_b\Delta_c(D_{2,2}D_{3,1}-D_{3,1}D_{1,2}-D_{2,2}D_{1,1}\\+D_{2,1}D_{1,2}+D_{3,2}D_{1,1}-D_{3,2}D_{2,1}),
% \end{multline*}
% where $\Delta_a=(-D_{2,1}D_{1,2}+D_{2,2}D_{1,1})$, $\Delta_b=(D_{2,2}D_{3,1}-D_{3,2}D_{2,1})$, and $\Delta_c=(D_{3,1}D_{1,2}-D_{3,2}D_{1,1})$. The determinant can become zero for different combination of values of the service demands, for instance, $D_{2,2}=D_{3,2}D_{2,1}/D_{3,1}$ gives $\Delta_b=0$, $\det(\comom{C}_{\comom{H}}^{1})=0$, and thus by (\ref{EQU:cblock}) this implies $\det({\bf \comom{A}}(\vec N))=0$. %A first technical solution consists in
% {\em Singularity Condition 1: Sparsity of Service Demands}.
% If some demands $D_{k,r}$ are equal to zero, then columns of all zeros appear in $\comom{\bf {A}}(\vec N)$ and the matrix becomes singular. A column of zeros  in $\comom{\bf {A}}(\vec N)$ implies that in (\ref{EQU:comom0}) the basis ${\Lambda}(\vec N)$ does not depend on all the normalizing constants in ${\Lambda}(\vec N-1_R)$. This case can be addressed by removing the normalizing constants associated to the zero columns both from all bases, which leaves an overdetermined system in (\ref{EQU:comom0}). Note that the solution methods discussed in Section \ref{SEC:solution} generalize also to the overdetermined case \cite{}.
%\begin{example}{\em
%\label{EXA:singzero}
% Consider a small model with the following parameters: $M=2$, $R=2$, $D_{1,1}=9$, $D_{1,2}=5$,  $D_{2,1}=0$,  $D_{2,2}=3$, $Z_{1}=6$, $Z_{2}=0$, and $\vec N=(8,1)$. The matrix ${\bf \comom{A}}(\vec N)$ and the basis $\Lambda(\vec N)$ are in this case:
% \[
% \small
% \renewcommand{\arraystretch}{1.1}
% \arraycolsep 0.1em
% \footnotesize
% {\bf \comom{A}}(\vec N)\Lambda(\vec N)=
% \begin{pmatrix}
%     -1 &~    1    &~ \cdot&~     \cdot   & -9 &~    \cdot    &~ \cdot &~    \cdot  &~   \cdot~\\
%     -1 &~    \cdot   &~  1 &~    \cdot   &~  \cdot &~    \cdot  &~   \cdot &~    \cdot  &~   \cdot~\\
%  ~    8 &~    \cdot    &~ \cdot  &  -6   & -9   &~  \cdot   &~  \cdot   &~  \cdot  &~   \cdot~\\
% ~     1  &~   \cdot   &~  \cdot  &~   \cdot    &~ \cdot  &~   \cdot &~    \cdot  &~   \cdot &~    \cdot~\\
% ~     \cdot   &~  \cdot  &~   \cdot   & -1    &~ 1 &~    \cdot &~    \cdot  &  -9  &~   \cdot~\\
% ~     \cdot   &~  \cdot   &~  \cdot &   -1   &~  \cdot  &~   1&~     \cdot  &~   \cdot &~    \cdot~\\
%   ~   \cdot   &~  \cdot  &~   \cdot &~    7 &~    \cdot  &~   \cdot  &  -6  &  -9  &~   \cdot~\\
%   ~   \cdot    &~ \cdot &~    \cdot  &~   1&~     \cdot   &~  \cdot &~    \cdot  &~   \cdot &~    \cdot~\\
%   ~   \cdot    &~ \cdot&~     \cdot   &~  \cdot&~     \cdot    &~ \cdot&~     1   &~  \cdot &~    \cdot~
% \end{pmatrix}
% %\Lambda(\vec N)=
% \begin{pmatrix}
% \small
% \renewcommand{\arraystretch}{1.1}
% \arraycolsep 0.1em
% G(\vec N)\\
% G(1_1,\vec N)\\
% G(1_2,\vec N)\\
% G(\vec N-\vec 1_1)\\
% G(1_1,\vec N-\vec 1_1)\\
% G(1_2,\vec N-\vec 1_1)\\
% G(\vec N-2\cdot \vec 1_1)\\
% G(1_1,\vec N-2\cdot \vec 1_1)\\
% G(1_2,\vec N-2\cdot \vec 1_1)\\
% \end{pmatrix}
% \]
% where $\cdot$ denotes zero entries. ${\bf \comom{A}}(\vec N)$ is singular being the last column associated to $G(1_2,\vec N-2\cdot \vec 1_1)$ of all zeros. Removing $G(1_2,\vec N-2\cdot \vec 1_1)$ from the basis $\Lambda(\vec N)$ provides an over-determined ${\bf \comom{A}}(\vec N)$ that has full rank and allows the recursive solution of the linear system (\ref{EQU:comom0}) as
% \begin{equation}
% {\Lambda}(\vec N)=[\comom{\bf A}^T(\vec N)\comom{\bf A}(\vec N)]^{-1}\comom{\bf A}^T(\vec N)\comom{\bf B}(\vec N){\Lambda}(\vec N-1_R),
% \end{equation}
% where we have pre-multiplied both sides by the transpose $\comom{\bf A}^T(\vec N)$ to make the coefficient matrix square.}
% \end{example}

% \subsubsection{Singularity Due to Identical Demands}
% ${\bf \comom{A}}(\vec N)$ also becomes singular if at least two service demands $D_{k_1,r}$ and $D_{k_2,r}$ of two queues $k_1$ and $k_2$ are identical for some classes $1\leq r\leq R-1$. This condition induces linear dependence between CEs and PCs sharing same constants of ${\lambda}(\vec 0) \subseteq {\Lambda}(\vec N)$. It is interesting to observe that also other methods such as RGF \cite{HarC02,HarL04} are complicated by the same condition, which suggests that the structure of the normalizing constant is significantly affected by this degeneration. We propose four strategies to address the problem:
% \begin{enumerate}
%     \item If the identical $D_{k,r}$'s all belong to a same class $r$, it is sufficient to rename class $r$ to class $R$ to prevent singularity.
%     \item If the identical $D_{k,r}$'s belong to different classes, e.g., the first $R^\prime$ classes, exact solutions may be still obtained by performing a MVA-like recursion on the first $R^\prime$ classes, while preserving the structure of the CoMoM recursion on the remaining $R-R^\prime$. This increases the computational costs by a factor $\prod_{r=1}^{R^\prime} (N_r+1)$.
%     \item if $m_k\leq M$ queues have {\em all} demands identical, i.e., the model includes $m_k-1$ replicas of queue $k$, it is sufficient to treat the model as composed by $M-m_k+1$ (distinct) queues and use in place of the PCs the more general relation
% \begin{equation}
% \label{EQU:genpc}
% \textstyle N_rG(\vec N)=Z_rG(\vec N-\vec 1_r)+\sum_{k=1}^M m_k D_{k,r}G(\vec 1_k,\vec N-\vec 1_r),
% \end{equation}
% for $1\leq r\leq R$, which compared to (\ref{EQU:pc}) accounts for the multiplicities $m_k$.
% The last formula is easily proved from Little's Law applied to the case of networks with queue replicas \cite{Cas06b}.
%
% \end{enumerate}

\section{Efficiency Analysis}
\label{SEC:efficiencyanalysis}
In this section we illustrate the increased efficiency of CoMoM with respect to MVA \cite{ReiL80}, multiclass Convolution \cite{ReiK75b}, RECAL \cite{ConG86}, and RGF\footnote{On models with $Z_r>0$ for some class $r$, RGF is run with $Z_r=0$ since no theoretical formulas are available to cope with think times in RGF.} \cite{HarC02}. We begin by two case studies of real multi-tier applications. In Case Study~1, we examine a model with six classes of requests and provide evidence that CoMoM is much more scalable than existing methods and can evaluate networks that are prohibitive with existing solution techniques. In Case Study~2, instead, we focus on a smaller case of model with three classes where MVA performs well and we prove that CoMoM can be valuable also on these small models. In Case Study~3, we show an example of model where the Wiedemann algorithm is much more scalable than Gaussian elimination. Finally, in Section \ref{SEC:parametric}, we explore the general scalability of CoMoM and compare it with previous work as the model parameters increase.

% \subsection{Computational Cost}
% \label{SEC:efficiency}
% In this section we give minimal complexity formulas for CoMoM time and space requirements. Focusing on the efficient implementation using fine decomposition, the computational costs of CoMoM are determined by three components:
% \begin{enumerate}
% \item at the beginning of the evaluation of a new class, the matrix $\comom{\bf A}_{1,1}$ must be inverted or factorized, with the latter choice being generally better in terms of sparsity. The LU factorization is required only for the diagonal micro-blocks and requires cubic time and quadratic space with respect to the micro-block orders. Therefore, it is of the order of ${M \choose h}^3M^3$ operations and ${M \choose h}^2M^2$ space.
% \item For each of the $N$ processed populations, we have to account the cost of LU back-substitution of the micro-blocks for evaluating (\ref{EQU:btfCoMoM}). This is less expensive than LU factorization and grows quadratically in the micro-block orders as ${M \choose h}^2M^2$; no significant storage overheads are required by LU back-substitution.
% \item Finally, we have to account that all operations are performed in exact linear algebra. Several methods for exactly representing numbers of arbitrary size exist, e.g., modulo representations. Here we focus on multiprecision algebra\footnote{Modulo arithmetic is often preferred to multiprecision algebra because it allows the parallel computation of linear systems. However, since the normalizing constant $G(\vec N)$ can reach astronomical values in large-scale multiclass models, even the prime numbers that define the modulo of an operation should be expressed as multiprecision numbers which significantly complicates implementations. For this reason, we believe that a simpler implementation of CoMoM in multiprecision algebra is more practical.}, where the exact linear solver carries on in computations all significant digits, considering the involved quantities as rational numbers. The computational overheads of multiprecision operations grow in complexity as $O^\sim(n)$, being $n$ the number of digits of the operands \cite{SchS71} and $O^\sim(n)$ the ``soft-Oh'' notation that hides polylogarithmic factors of the type $\log N$,  $\log \log N$, $\log \log \log N$, etc. For simplicity of notation, we here use the standard $O(\cdot)$ notation and implicitly assume to ignore all polylogarithmic factors except for the leading factor $\log N$. Since normalizing constants digits grow at most as $O(N\log N)$ \cite{Cas06b}, this implies that the computational overheads of exact algebra grow approximately as $O(N\log N)$ both in time and in storage.
% \end{enumerate}
% Combining the above observations, we find that the time requirements of CoMoM using the fine decomposition are approximately of the order of $\sum_{h=0}^{\comom{H}} {R-1\choose h} {M \choose h}^3M^3+ N^2\log N\sum_{h=0}^{\comom{H}} {R-1\choose h} {M \choose h}^2M^2.$ The analysis of memory requirements is similar. The storage costs are due to the exact multiprecision entries of the matrices $\comom{\bf A}_{i,j}$ and  $\comom{\bf B}$ grow approximately as $N\log N\sum_{h=0}^{\comom{H}} {R-1\choose h} {M \choose h}^2M^2,$ where $N\log N$ accounts for the multiprecision overheads, while the right-hand side is the maximum storage occupation of the micro-blocks. In fact, both $\comom{\bf B}$ and the non-diagonal blocks of $\comom{\bf A}_{i,j}$ are extremely sparse and their contribution to the overall memory requirements can be ignored\footnote{It is also possible to observe that since our storage complexity formula assumes the micro-blocks as dense matrices, it is already overestimating the actual memory complexity, since in general the number of zero elements in the micro-blocks are more than the number of non-zeros in $\comom{\bf B}$ and in the non-diagonal blocks.}.
% Recalling that existing algorithms such as MVA \cite{ReiL80} or {\sc Recal} \cite{ConG86} grow as $O(N^R)$ and $O(N^M)$ with the population size in both time and space, it is immediate to conclude that the proposed algorithm is much more efficient, being approximately log-quadratic in time and log-linear in space with respect to $N$. We build intuition on the entity of the computational gains of CoMoM in the next subsections using a case study and parametric plots.
\begin{figure}[t]
\centering
  % Requires \usepackage{graphicx}
{\includegraphics[width=0.8\columnwidth]{fig3}}
  \caption{{\bf Multi-tier Application Case Study}. Queueing network model of a real J2EE e-business system \cite{KouB03}.}
  \label{FIG:j2ee}
\end{figure}


\subsection{Case Study 1: J2EE e-Business Model}
We illustrate CoMoM on the queueing network model presented in \cite{KouB03} of a real J2EE e-business system, see Figure~\ref{FIG:j2ee}. The model captures the performance of a real software system in execution on a multi-tier architecture composed by a cluster of nine application servers (AS) and a dual-processor database server (DB). Queues indexed $1-9$ are application servers, queues $10-11$ represent database processors, queue $12$ models HTTP communication overheads. The network does not include communication loops between the AS and the DB because their effect on performance is already accounted within the mean number of visits that define the service demands $D_{k,r}$ \cite{KouB03,LazZGS84}. All queues use either processor sharing (queues $1-11$) or first-come first-served  scheduling (queue $12$), with the exception of the delay server where jobs wait in think state (no queueing) before re-entering the network. The system processes $R=5$ classes of business workloads:
\begin{itemize}
\item {\sf NewOrder}, enters a new order in the system (class C1);
\item {\sf ChangeOrder}, changes an existing order (class C2);
\item  {\sf OrderStatus}, reads the status of an existing order (class C3);
\item  {\sf CustStatus},  lists orders of a given customer (class C4);
\item  {\sf WorkOrder}, enters a new manufacturing order (class C5);
\end{itemize}
 The service demands found in \cite{KouB03} for the five classes are given in Table~\ref{TAB:casestudy}.  We set the mean think times to $Z_1=Z_2=Z_3=Z_4=2s$ and $Z_5=3s$, respectively.
\begin{table}[ht]
\centering
\caption{Service demands  $D_{k,r}$ [ms] of the J2EE application model}
\label{TAB:casestudy}
\begin{tabular}{|c|cc ccc|}
\hline
\emph{\textrm{Queues} vs {\rm Classes}} & \emph{C1} &\emph{C2} &\emph{C3} &\emph{C4} &\emph{C5} \\
   \hline
1--9 &12.98 & 13.64 & 2.64 & 2.54 & 24.22 \\
10--11 & 5.32 & 5.18 & 1.24 & 1.04 & 17.07 \\
12& 1.12 & 1.27 & 0.58 & 0.03 & 1.68 \\
\hline
\end{tabular}
\end{table}

All population are kept of identical size while increasing $N$, i.e., $N_1=N_2=\ldots=N_{R}=N/R$. Because some queues are replicated, we use in CoMoM the generalized PC (\ref{EQU:genpc}). The comparison of CoMoM with the computational costs of MVA \cite{ReiL80}, multiclass Convolution \cite{ReiK75b}, RECAL \cite{ConG86} and RGF \cite{HarL04} is given in Table~\ref{TAB:j2eecomparison}. The table reports time and space requirements for the different methods. Results are rounded up to the closest larger integer. The notation ``limit'' indicates that the algorithm has exceed either time or space requirements which are set to 1GB of RAM and $10$ minutes on a Intel Core Duo $2\times 1.60$ GHz, respectively; ``n/a'' indicates that the value cannot be measured because of the infeasible requirement of another physical resource (time or memory). Numerical stabilization of the normalizing constant for existing methods is done by scaling of the service demands\cite{Lam82,ConG86}. We have experimented both with Gaussian elimination based on LU back-substitution and with the Wiedemann algorithm. On this example, we report results from the Gaussian elimination implementation which has been found at least one order of magnitude faster than the Wiedemann algorithm in all experiments, e.g., for $N=100$ the Wiedemann algorithm takes $39$s instead of the $1$s of Gaussian elimination. The results in Table~\ref{TAB:j2eecomparison} indicate that CoMoM is the only method capable of solving models for all ranges of populations considered in the experiments. For larger populations values that the ones considered here, the gains of CoMoM are even larger, e.g., for $N=5000$ theoretical formulas indicate that {the MVA time costs would be about five orders of magnitude larger than CoMoM and the memory occupation would be eight orders of magnitude larger}. We remark that the much larger requirements of RECAL with respect to the other methods derive from its inefficiency in processing networks with queue replicas\cite{HarL04}. However, also using the initialization of RECAL proposed in \cite{Cas06b} for models with queue replicas, RECAL remains less scalable than RGF.
%Summarizing, the case study indicates that CoMoM scales effectively on models that find application in the evaluation of real software systems and which are prohibitive for existing methods. We also remark that even bigger computational gains can be obtained as the model complexity grows. %A parametric analysis of CoMoM efficiency is given in the next subsection.

\begin{table}[t]
\centering
\renewcommand{\arraystretch}{1.1}
\caption{J2EE System Case Study}
\label{TAB:j2eecomparison}
\begin{tabular}{|c|c|c|c|c|c|}
\hline
 %&  &  &  &  &  \\
\emph{$N$ vs {Time}} & \emph{MVA} & \emph{Conv.} & \emph{RECAL} & \emph{RGF} & \emph{CoMoM}/LU\\%/Wiedemann) \\
\hline
$10$ & $1$s & $1$s & $1$s & $1$s  & $1$s\\%/$<1$s \\
$50$ & $1$s & $1$s & limit & $1$s  & $1$s\\%/$10$s \\
$100$ & $2$s & $1$s & limit & $5$s  & $2$s\\%/$39$s \\
$500$ & n/a & n/a & limit & limit  & $52$s\\%/$864$s \\
$1000$ & n/a & n/a & limit & limit  & $327$s\\%/$(0)$s \\
%$5000$ & n/a & n/a & limit & limit  & $871$s\\%/$(0)$s \\
\hline
\emph{$N$ vs {Space}} & \emph{MVA} & \emph{Conv.} & \emph{RECAL} & \emph{RGF} & \emph{CoMoM}/LU \\
\hline
$10$ & $1$MB & $1$MB & $1$MB & $1$MB  & $1$MB\\%/$2$MB \\
$50$ & $7$MB & $2$MB & n/a & $1$MB  & $2$MB\\%/$2$MB \\
$100$ & $140$MB & $43$MB & n/a & $1$MB  & $5$MB\\%/$2$MB \\
$500$ & limit & limit & n/a & n/a  & $7$MB\\%/$0$ \\
$1000$ & limit & limit & n/a & n/a  & $9$MB \\
%$5000$ & limit & limit & n/a & n/a  & $24$MB \\
%\hline
%\multicolumn{6}{|l|}{{\em Remarks:} The best results of each row are underlined. Results are}\\
%\multicolumn{6}{|l|}{rounded toward the closest bigger integer.}\\
\hline
\end{tabular}
\end{table}

%\begin{table}[ht!]
%\centering
%\renewcommand{\arraystretch}{1.1}
%\caption{J2EE System Case Study: Balanced Growth}
%\label{TAB:j2eecomparison}
%\begin{tabular}{|c|c|c|c|c|c|}
%\hline
%$N$ vs {Time}   & MVA & Conv. & {\sc Recal} & RGF  & CoMoM \\
%\hline
%$10$ & $\underline{10^4}$ & $\underline{10^4}$ & $10^7$ & $10^5$  & $10^6$ \\
%$50$ & $\underline{10^7}$ & $10^7$ & $(10^{14})$ & $10^7$ & $10^7$ \\
%$100$ & $10^8$ & $10^8$ & $(10^{17})$ & $10^9$ &  $\underline{10^7}$ \\
%$500$ & $10^{11}$ & $10^{11}$ & $(10^{25})$ & $10^{11}$  & $\underline{10^9}$ \\
%$1000$ &  $(10^{13})$ & $(10^{13})$ & $(10^{29})$ & $(10^{13})$ & $\underline{10^{10}}$ \\
%$5000$ &  $(10^{16})$ & $(10^{16})$ & $(10^{37})$ & $(10^{15})$ & $\underline{10^{11}}$ \\
%$10000$ & $(10^{18})$ & $(10^{18})$ & $(10^{41})$ & $(10^{17})$ & $\underline{10^{12}}$ \\
%\hline
%$N$ vs {Space}   & MVA & Conv. & {\sc Recal} & RGF & CoMoM \\
%\hline
%$10$ & $10^3$ & $10^3$ & $10^3$ & $\underline{10^2}$ & $10^5$ \\
%$50$ & $10^6$ & $10^6$ & $(10^5)$ & $\underline{10^4}$ & $10^6$ \\
%$100$ & $10^8$ & $10^8$ & $(10^7)$ & $\underline{10^5}$ & $10^6$ \\
%$500$ & $(10^{11})$ & $(10^{11})$ & $(10^9)$ & $\underline{10^7}$ & $\underline{10^7}$ \\
%$1000$ &  $(10^{13})$ & $(10^{13})$ & $(10^{11})$ & $(10^{8})$  & $\underline{10^{7}}$ \\
%$5000$ &  $(10^{16})$ & $(10^{16})$ & $(10^{13})$ & $(10^{10})$  & $\underline{10^{8}}$ \\
%$10000$ & $(10^{18})$ & $(10^{18})$ & $(10^{15})$ & $(10^{11})$  & $\underline{10^{8}}$ \\
%\hline
%\end{tabular}
%\end{table}

\subsection{Case Study 2: Two-Tier PeopleSoft Application Model}

We illustrate a case study of multi-tier architecture where CoMoM shows increased effectiveness with respect to the MVA algorithm also on models with a small number of classes ($R=3$). Differently from Case Study 1, we here consider the basic implementation of (\ref{EQU:comom0}) without BTFs. Our aim is to explore a small model where the simplest implementation of CoMoM starts to be more effective than MVA.

We consider a queueing network similar to the two-tier application model in \cite{ChuAT97} of an Oracle PeopleSoft client-server system. Client applications access a remote database server and are modeled in the queueing network as a delay server. The database server is instead modeled by two queues representing its CPU and disk drive. Because we focus on constant-rate servers, we represent the network bandwidth contention by a processor-sharing queue instead of the load-dependent station in \cite{ChuAT97}. The related queueing network model is shown in Figure~\ref{FIG:peoplesoft}. We use the benchmarking characterization and measurements in \cite{ChuAT97} to parameterize the queueing network. The benchmark has three classes of workloads, {\sf Stress Insert}, {\sf Stress Update}, and {\sf Stress Delete} which differ for the type of SQL operations in the workload \cite{ChuAT97}. The service demands for the different classes and resources given in \cite{ChuAT97} are reported in Table~\ref{TAB:casestudy2}. We set the think times at the delay server to $Z_1=Z_2=Z_3=5s$.
\begin{table}[ht]
\centering
\caption{Service demands  $D_{k,r}$ [ms] of the Oracle PeopleSoft model}
\label{TAB:casestudy2}
\begin{tabular}{|c|ccc|}
\hline
\emph{\textrm{Queues} vs {\rm Classes}} & C1 &C2 &C3 \\
   \hline
Cpu & 12.21 & 21.73 & 21.96  \\
Disk & 26.98 & 31.94 & 29.24  \\
Net & 16.12 & 26.46 & 28.51 \\
\hline
\end{tabular}
\end{table}

We have performed a comparison of CoMoM and MVA using an approach similar to Case Study 1, but without the use of BTFs. Also in this case, the implementation based on LU back-substitution was significantly faster than the Wiedemann algorithm (at least two orders of magnitude) because of the small size of the model under consideration. The results, shown in Table~\ref{TAB:peoplesoftcomparison}, further confirm the increased scalability of CoMoM with respect to MVA. MVA becomes memory inefficient as soon as the total population grows beyond $900$ jobs, whereas the memory occupation of CoMoM in all experiments is always less than $4$MB. The comparison of computational times is instead favorable to the MVA, thanks to the small number of classes which make the MVA recursion not too expensive. Nevertheless, on this small model, the requirements of CoMoM remain good also without BTFs. This indicates that, also using the simplest possible implementation of (\ref{EQU:comom0}), one may obtain a solution algorithm that is much more memory efficient than MVA and with similar running times.

%Summairizing, the results of this section suggest that the CoMoM technique in the basic implementation (\ref{EQU:comom0}) can be valuable also on very small models when these are populated by several hundreds requests.

\begin{figure}[t]
\centering
  % Requires \usepackage{graphicx}
{\includegraphics[width=0.9\columnwidth]{fig4}}
  \caption{{\bf Oracle PeopleSoft Application Case Study}. Queueing network model of a two-tier PeopleSoft client-server system \cite{ChuAT97}.}
  \label{FIG:peoplesoft}
\end{figure}

\subsection{Case Study 3: Applicability of the Wiedemann Algorithm}
\label{SEC:wiedemann}
In the previous case studies, we have reported that Gaussian elimination based on LU factorization and back-substitution performed better than the Wiedemann Algorithm recommended in Section \ref{SEC:linalg} for the evaluation of large-scale models. In this section, we show cases of linear systems (\ref{EQU:comom0}) where the Wiedemann algorithm outperforms Gaussian elimination.

We consider a model with $M=4$ queues and $R=8$ classes and synthetic service demands $D_{k,r}=k^r$. We evaluate the solution of a single recursive step of (\ref{EQU:comom0}) with Gaussian elimination and with the Wiedemann algorithm. From Table 2, it is possible to see that ${\bf A}(\vec N)$ in this case has order $1650$ with $9494$ non-zeros. For example, we have performed experiments where Gaussian elimination requires $17$s of CPU time and $244$MB of memory to complete a single recursive step when processing class $8$. On the same machine, the Wiedemann algorithm performs the same recursive step in $9$s and the memory occupation is dramatically lowered to $4$MB. This is a consequence of the fact that the Wiedemann algorithm is an iterative methods which preserves the sparsity of ${\bf A}(\vec N)$.

We have also performed the same experiment for $M=4$ and $R=10$ and we have observed that Gaussian elimination becomes infeasible due to the high computational requirements that would require approximately $220$s and $1.2$GB for the solution. The Wiedemann algorithm is instead able to solve the recursive step in $69s$ with only $7$MB of memory occupation. Note that these values are relative to the evaluation of  (\ref{EQU:comom0}) without BTF; whenever BTF is used, the advantages of the Wiedemann algorithm become evident only when the order of the micro-blocks is of the order of thousands, see Table \ref{TAB:Areduction}.

\begin{table}[t]
\centering
\renewcommand{\arraystretch}{1.1}
\caption{PeopleSoft Application Case Study}
\label{TAB:peoplesoftcomparison}
\begin{tabular}{|c|c|c|c|c|}
\hline
 & \multicolumn{2}{c|}{\em Time [s] - no BTF} & \multicolumn{2}{c|}{\em Space [MB] - no BTF} \\
\hline
\emph{$N$}& \emph{MVA} & \emph{CoMoM}/LU & \emph{MVA} & \emph{CoMoM}/LU \\
\hline
% $100$ & $<1$s & $<1$s/$<1$s & $2$MB & $3$MB/$2$MB \\
% $300$ & $<1$s & $2$s/$71$s & $38$MB & $3$MB/$2$MB \\
% $600$ & $5.5$s & $11$s/$479$s & $280$MB & $4$MB/$2$MB \\
% $900$ & $19.9$s & $30$s/$1793$s & $905$MB & $4$MB/$3$MB \\
% $1200$ & Infeas. & $62$s/$4983$s & Infeas. ($>2$GB) & $4$MB/$3$MB \\
$100$ & $1$s & $1$s & $2$MB & $3$MB \\
$300$ & $1$s & $5$s & $38$MB & $3$MB \\
$600$ & $3$s & $27$s & $280$MB & $3$MB \\
$900$ & $11$s & $82$s & $938$MB & $4$MB \\
$1200$ & n/a & $191$s & limit & $4$MB \\
%\hline
%\multicolumn{5}{|l|}{{\em Remarks:} Results between parenthesis are estimates for infeasible models.}\\
%\multicolumn{5}{|l|}{{\em Remarks:} The best results of each row are underlined. Results are}\\
%\multicolumn{5}{|l|}{rounded toward the closest bigger integer.}\\
\hline
\end{tabular}
\end{table}
\subsection{Parametric Analysis of Computational Costs}
\label{SEC:parametric}
We conclude the section by showing with a parametric analysis that the computational requirements of CoMoM scale efficiently with the number of service classes. The requirements of CoMoM are compared to those of MVA and RGF in Figure~\ref{FIG:parametric} for increasing values of the number of classes $R$. In the experiments, we vary the number of classes in $R=1,2,\ldots,12$ and set $M=3$, $N=1000$, $N_r=N/R$, $D_{max}=1$ and $Z_r=0$, $1\leq r\leq R$. Experiments with different populations or number of queues give results that are usually even more favorable to CoMoM. In the evaluation of computational requirements, we do not consider the Wiedemann algorithm because this is a probabilistic algorithm and hence its computational requirements vary significantly depending on the actual value of the service demands $D_{k,r}$. We instead focus on the CoMoM implementation based on Gaussian elimination (LU factorization and back-substitution), which provides a worst-case estimate of CoMoM behavior in terms of memory. In order to show the expect behavior of the methods also on models that are infeasible, estimates are obtained from theoretical formulas.

% In order to simplify comparison of CoMoM with results reported in the literature, which are expressed in terms of number of operations and storage units, we use these indicators in Figure~\ref{FIG:parametric}.  For the hardware architecture used, Pentium 4 3.60 GHz, the number of operations performed in one second is approximately $1\cdot 10^8$. For the storage units, $10^{9}$ is approximately $2$GB of RAM. Here we assume that all classes have identical population $N_r$.

The figures show that CoMoM has typically the lowest requirements among the considered methods as $R$ grows beyond $3$ or $4$ classes. MVA and RGF can be better only for models with a very small number of classes. In general, CoMoM becomes inefficient only if both $M$ and $R$ grow simultaneously, because the size of the blocks (\ref{EQU:cblock}) mostly depends on $\comom{H}=\min\{M,R-1\}$, but this inefficiency is present also in all existing algorithms and remains an open issue of multiclass models.

The analysis of storage requirements indicates that RGF has minimal memory costs. However, CoMoM is much faster than RGF while keeping good memory requirements that are also dramatically lower than MVA memory occupation. We have found with additional analyses that, as the number of queue $M$ grows, the memory requirements of CoMoM and RGF become similar, while the time costs remain much more favorable to CoMoM, we point to \cite{TR} for details. We also remark that, for the models where CoMoM has the largest memory requirements, one may simply use the implementation based on Wiedemann algorithm that has even negligible storage costs, see the formulas in Section \ref{SEC:solution} for memory occupation estimates in this case.

\section{Evaluation of State-Dependent Cost Indexes}
\label{SEC:statedepcosts}
In this section, we define a generalization of CoMoM, called Probabilistic CoMoM (Pro-CoMoM), that can efficiently compute marginal probabilities in multiclass models. The marginal probability of observing $n$ jobs in queue $k$ is defined as~\cite{ReiK75b}
\begin{equation}
\textstyle \Pr_k(n\,|\, \vec N)=\sum_{\vec n_k: n_{k}=n} F_k(\vec n_k) {G(-\vec 1_k,\vec N-\vec n_k)}/{G(\vec N)},
\end{equation}
where $F_k(\vec n_k)=n_k!\prod_{r=1}^R {D_{k,r}^{n_{k,r}}}/{n_{k,r}!}$.
Pro-CoMoM computes marginal probabilities by first obtaining the un-normalized values
$$\widetilde{\Pr}_k(n\,|\, \vec N)=\Pr_k(n\,|\, \vec N)G(\vec N)$$ and then scaling back these terms into probabilities by imposing $$\textstyle G(\vec N)=\sum_{n=0}^N\widetilde{\Pr}_k(n\,|\, \vec N)$$ that follows from the condition $\sum_{n=0}^N\Pr_k(n\,|\, \vec N)=1$.

We obtain Pro-CoMoM in two steps. First, we generalize CEs and PCs to the computation of marginal queue-length probabilities. Then, we combine these equations into linear systems to obtain a recursive computational scheme similar to CoMoM.
\begin{figure}[t]
\centering
  % Requires \usepackage{graphicx}
  %\subfigure[$M=3$, $N_r=20\; \forall r$]{\includegraphics[width=0.24 \textwidth]{compareM3N20}}
  \subfigure[]{\includegraphics[width=0.24 \textwidth]{fig5a}}
  %\subfigure[$M=3$, $N_r=100\; \forall r$]{\includegraphics[width=0.24 \textwidth]{compareM3N100}}
  %\subfigure[$M=5$, $N_r=100\; \forall r$]{\includegraphics[width=0.24 \textwidth]{compareM5N100}}\\
  %\subfigure[$M=3$, $N_r=20\; \forall r$]{\includegraphics[width=0.24 \textwidth]{compareM3N20space}}
  \subfigure[]{\includegraphics[width=0.24 \textwidth]{fig5b}}
  %\subfigure[$M=3$, $N_r=100\; \forall r$]{\includegraphics[width=0.24 \textwidth]{compareM3N100space}}
  %\subfigure[$M=5$, $N_r=100\; \forall r$]{\includegraphics[width=0.24 \textwidth]{compareM5N100space}}\\
  \caption{{\bf Parametric Analysis of CoMoM Requirements}. Parametric comparison of computational requirements for different values of the number of classes. On the considered model: $M=3$, $N=1000$, $N_r=N/R$, $1\leq r \leq R$. Results for models with different number of queues or populations are qualitatively similar.}% Estimates refer to an architecture that can process $10^8$ operations per second.}
  \label{FIG:parametric}
\end{figure}
{\em Generalization of CEs and PCs}. The fundamental prerequisite for the development of a recursive scheme for marginal queue-lengths similar to CoMoM is the existence of equations similar to the CEs and PCs and which relate marginal probabilities instead of normalizing constants. Theorem~\ref{THM:Pmarg} investigates this generalization.
\begin{theorem}
\label{THM:Pmarg}
{\em The un-normalized probabilities $\widetilde{\Pr}_k(n\,|\, \vec N)$, $1\leq n\leq N$, satisfy the probabilistic convolution expressions (PCE)
\begin{align}
\label{EQU:PCE}
\textstyle \widetilde{P}_k(n\,|\,\vec 1_j, \vec N)=&\textstyle \widetilde{P}_k(n\,|\,\vec N)+\sum_{r=1}^R D_{j,r}\widetilde{P}_k(n\,|\, \vec 1_j, \vec N-\vec 1_r),
\end{align}
for all $1\leq n\leq N$, $1\leq k\leq M$, $1\leq j\leq M$, $j\neq k$, and the probabilistic population constraints (PPC)
\begin{multline}
\label{EQU:PPC}
\textstyle N_r\widetilde{P}_k(n\,|\, \vec N)=Z_r\widetilde{P}_k(n\,|\, \vec N-\vec 1_r)+nD_{k,r}\widetilde{P}_k(n-1\,|\, \vec N-\vec 1_r)\\+\textstyle \sum_{j\neq k}D_{j,r}\widetilde{P}_k(n\,|\, \vec 1_j,\vec N-\vec 1_r),
\end{multline}
for arbitrary $1\leq r\leq R$, where $\widetilde{P}_k(n\,|\, \vec 1_j,\vec N)$ and $\widetilde{P}_k(n-1\,|\, \vec 1_j,\vec N-\vec 1_r)$ refer to  models with a replica of queue $j$ more.
}
\end{theorem}
\begin{proof}
We use the notation, e.g., $G(\vec 1_j-\vec 1_k,\vec N)$, to describe the normalizing constant of a model obtained from the original network by adding a replica of queue $j$ and removing queue $k$.
Formula (\ref{EQU:PCE}) follows by writing
\begin{equation}
\textstyle \widetilde{P}_k(n\,|\,\vec 1_j, \vec N)=\sum_{\vec n_k:n_k=n} F_k(\vec n_k)G(\vec 1_j-\vec 1_k,\vec N-\vec n_k), \nonumber
\end{equation}
and applying to $G(\vec 1_j-\vec 1_k,\vec N-\vec n_k)$ the CE for queue $j$, i.e.,
\begin{multline}
\textstyle G(\vec 1_j-\vec 1_k,\vec N-\vec n_k)=G(-\vec 1_k, \vec N-\vec n_k)\\\textstyle +\sum_{r=1}^R D_{j,r}G(\vec 1_j-\vec 1_k,\vec N-\vec 1_r-\vec n_k). \nonumber
\end{multline}
The resulting convolution with $F_k$ gives immediately (\ref{EQU:PCE}).

Instead, the derivation of (\ref{EQU:PPC}) starts by observing that
\begin{multline}
\textstyle  \sum_{\vec n_k:n_k=n} n_{k,r}F_k(\vec n_k)G(-\vec 1_k,\vec N-\vec n_k)=nD_{k,r}\widetilde{P}(n-1\,|\,\vec N-\vec 1_r). \nonumber
\end{multline}
which implies the following equivalence
\begin{multline}
\label{TMP:1}
\textstyle  \sum_{\vec n_k:n_k=n} (N_r - n_{k,r})F_k(\vec n_k)G(-\vec 1_k,\vec N-\vec n_k)\\
\textstyle =N_r\widetilde{P}_k(n\,|\, \vec N)-nD_{k,r}\widetilde{P}_k(n-1\,|\, \vec N-\vec 1_r).
\end{multline}
Noting that $G(-\vec 1_k,\vec N-\vec n_k)$ in the left-hand side of (\ref{TMP:1}) can be expanded by (\ref{EQU:pc}) as
$$\textstyle Z_{r}G(-\vec 1_k,\vec N-\vec 1_r-\vec n_k)+\sum_{j\neq k} D_{j,r}G(\vec 1_j-\vec 1_k,\vec N-\vec 1_r-\vec n_k),$$
which accounts for the fact that queue $k$ is not anymore in the network, the proof of (\ref{EQU:PPC}) follows by applying the last expansion to the left-hand side of (\ref{TMP:1}) and observing that the resulting convolutions over $\vec n_k$ immediately give (\ref{EQU:PPC}).
\end{proof}

%\begin{figure}[ht!]
%\small
%\begin{center}
%\framebox
%{\begin{minipage}[h]{\columnwidth}
%\begin{tabbing}
%9\=99\=99\=99\=99\=99\=99\=999999999999999999999999999999999999999999999999\=99999999999999999999999999999999999999\=\kill
%Set probabilities with negative populations in $\Pi_k(\vec 0)$ to zero \\
%Set probabilities with population $\vec 0$ in $\Pi_k(\vec 0)$ for $n_k=0$ to one \\
%Set probabilities with population $\vec 0$ in $\Pi_k(\vec 0)$ for $n_k>0$ to zero \\
%\> {\bf for } $r$ {\bf from} $1$ {\bf to} $R$\\
%\> \> {\bf for } $n_r$ {\bf from} $1$ {\bf to} $N_r$\\
%\> \> \> $\vec n=(N_1,N_2,\ldots,N_{r-1},n_r,\vec 0,\ldots,\vec 0)$\\
%\> \> \> Compute $\comom{\bf A}_k(\vec n)$, $\comom{\bf B}_k(\vec n)$, $\comom{\bf C}_k(\vec n)$, $\comom{\bf D}_k(\vec n)$\\
%\> \> \> {\bf for } $n$ {\bf from} $1$ {\bf to} $N_1+N_2+\ldots+N_{r-1}+n_r$\\
%\> \> \> \> $\Pi_k(n\,|\,\vec n)=\comom{\bf B}_k(n\,|\,\vec n)\Pi_k(n\,|\,\vec n-\vec 1_r)$\\
%\> \> \> \> $\Pi_k(n\,|\,\vec n)=\Pi_k(n\,|\,\vec n)+\comom{\bf C}_k(n\,|\,\vec n)\Pi_k(n-1\,|\,\vec n)$\\
%\> \> \> \> $\Pi_k(n\,|\,\vec n)=\Pi_k(n\,|\,\vec n)+\comom{\bf D}_k(n\,|\,\vec n)\Pi_k(n-1\,|\,\vec n-\vec 1_r)$\\
%\> \> \> \> $\Pi_k(n\,|\,\vec n)=\comom{\bf A}_k^{-1}(n\,|\,\vec n)\Pi_k(n\,|\,\vec n)$\\
%\> \> \> {\bf end for}\\
%\> \> {\bf end for}\\
%\> {\bf end for}\\
%\> $G(\vec N)=0$\\
%\> {\bf for } $n$ {\bf from} $1$ {\bf to} $N$\\
%\> \> $G(\vec N)=G(\vec N)+\Pi_k(n\,|\,\vec n)$\\
%\> {\bf end for}\\
%\> {\bf for } $n$ {\bf from} $1$ {\bf to} $N$\\
%\> \> Extract $\widetilde{P}_k(n\,|\,\vec N)$ from $\Pi_k(n\,|\,\vec N)$\\
%\> \> $P_k(n\,|\,\vec N)=\widetilde{P}_k(n\,|\,\vec N)/G(\vec N)$\\
%\> {\bf end for}\\
%\> Return $P_k(n\,|\,\vec N)$, $0\leq n\leq N$
%\end{tabbing}
%\end{minipage}}
%\end{center}
%\caption{{\bf Basic Pro-CoMoM implementation for computing $P_k(n\,|\,\vec N)$.}
%%Pseudo-code for the computation of $\comom{\bf A}(\vec n)$ and $\comom{\bf B}(\vec n)$ is reported in the final Appendix.
%Compared to CoMoM, Pro-CoMoM requires an additional inner loop for the evaluation of the marginal queue-length probabilities of queue $k$ and final loops for re-scaling the un-normalized values by $G(\vec N)$.}
%\label{ALG:procomom}
%\end{figure}
The PCEs and PPCs equations provide generalizations of the CEs and PCs that hold also outside their application in the Pro-CoMoM algorithm. In particular, one may solve any of the two equations in isolation by a recursive scheme similar to MVA/LBANC for the PCEs, or to RECAL for the PPCs, and obtain the value of the marginal probabilities $P_k(n|\vec N)$. While the computation of marginal probabilities with a recursion similar to the MVA has been already explored in the literature (load-dependent MVA, \cite{ReiL80}), a recursive evaluation of the PPCs would produce an algorithm with different computational trade-offs with respect to state-of-the-art methods, i.e., which grows efficiently with the number of classes $R$, but exponentially with the number of queues $M$. However, this appears of limited practical applicability due to the high cost of recursions similar to RECAL on models with more than a few tens of requests\footnote{Although we here focus on the extension of CoMoM, the development of PCEs and PPCs suggests that it may also be possible to derive also a generalization of the MoM algorithm in \cite{Cas06b} to compute marginal queue-length probabilities. However, this generalization of MoM lies outside the scope of this paper and it is left as future work.}.

{\em Pro-CoMoM Algorithm.} We now define the Pro-CoMoM algorithm by solving PCEs and PPCs into a recursive linear system of equations similar to (\ref{EQU:comom0}). We proceed similarly to the derivation of CoMoM by first defining a basis of marginal queue-length probabilities, and then describing a matrix equation for its recursive computation.
\begin{definition} {\em The {\em probabilistic basis} $\Pi_k(n\,|\,\vec N)$ for queue $k$  is the set of un-normalized marginal probabilities
\begin{equation}
\Pi_k(n\,|\,\vec N)=\pi_k(n,\vec 0) \bigcup_{1\leq j\leq M\,\wedge\,j\neq k} \pi_k(n,\vec 1_j),
\end{equation}
where $\pi_k(n,\vec v)=\{\widetilde{P}_k(n\,|\,\vec v,\vec N- \vec n) \,|\, \vec v\neq \vec 1_k \,\wedge\, \sum_{r=1}^{R-1} n_r \leq M \,\wedge\, n_R=0\}$ is the set of un-normalized marginal probabilities associated to models with
$\vec v$ queues replicas.
}
We remark that the values $\widetilde{P}_k(n-1\,|\, \vec 1_k, \vec N-\vec 1_r)$, $1\leq n \leq N$, $1\leq k\leq M$, are not used in the PPCs and thus are omitted from the basis $\Pi_k(n\,|\,\vec N)$. The Pro-CoMoM recursion is introduced below. Pseudo-code of the recursions and algorithms for generating the matrices used in the recursive step are reported in \cite{TR}. We also note that, since (\ref{EQU:PCE})-(\ref{EQU:PPC}) involve only marginal probabilities of the same queue $k$, the analysis of a queue marginals can be performed independently from the other queues. A complete evaluation of all marginal probabilities is then performed by $M$ successive runs of Pro-CoMoM for $k=1,2,\ldots,M$.
\begin{theorem} {\em The probabilistic basis $\Pi_k(n\,|\,\vec N)$ satisfies
\begin{multline*}
\comom{\bf A}_k(n\,|\,\vec N) \Pi_k(n\,|\,\vec N)=\comom{\bf B}_k(n\,|\,\vec N)\Pi_k(n\,|\,\vec N-\vec 1_r)\\+\comom{\bf C}_k(n\,|\,\vec N)\Pi_k(n-1\,|\,\vec N)+\comom{\bf D}_k(n\,|\,\vec N)\Pi_k(n-1\,|\,\vec N-\vec 1_r),
\end{multline*}
for all $0\leq n\leq N$ and $1\leq k\leq M$, where the four matrices have all order $M{M+R-1\choose M}$ and are defined by the coefficients of all equations (\ref{EQU:PCE})-(\ref{EQU:PPC}) relating $\Pi_k(n\,|\,\vec N)$, $\Pi_k(n-1\,|\,\vec N)$, $\Pi_k(n\,|\,\vec N-\vec 1_r)$, and $\Pi_k(n-1\,|\,\vec N-\vec 1_r)$.}
\end{theorem}
\begin{proof}Similarly to the approach in Theorem \ref{THM:CoMoM}, $\pi_k(0)$ is first computed using the PPC of class $R$. We prove the rest of the statement by induction on the set of known marginal probabilities. \\
\noindent {\em Case $n=0$}. In this case (\ref{EQU:PCE})-(\ref{EQU:PPC}) have the same identical structure of (\ref{EQU:ce})-(\ref{EQU:pc}) except for the missing marginal probabilities $\Pr_k(0\,|\,\vec 1_k,\vec N-\vec N)$ which reduce the number of unknowns in the set difference $\Pi_k(\vec N)/\pi_k(0)$ to $(M-1){M+R-1 \choose M}.$ Similarly to the proof of Theorem \ref{THM:CoMoM}, we can therefore formulate $(M+R-2){M+R-2 \choose M-1},$
where the first coefficient is now $(M+R-2)$ instead of $(M+R-1)$ because we can formulate only $M-1$ PCEs due to the missing marginal probabilities. By comparison of the number of unknowns with the number of rows, we find that the degrees of freedom are
$df=(M-1){M+R-1 \choose M}-(M+R-2){M+R-2 \choose M-1}={M+R-2 \choose M-1}-{M+R-1 \choose M}\leq 0$ equations, which make the linear system always square or over-determined.

\indent  {\em Inductive step}. We assume that, when processing the queue-length size $n$, the solution for $n-1$ is already available. In this case (\ref{EQU:PCE})-(\ref{EQU:PPC}) have again the same structure of (\ref{EQU:ce})-(\ref{EQU:pc}) since the terms depending on the marginal probabilities for queue-length $n-1$ are known. Therefore, the same considerations for $n=0$ apply readily to this case and complete the proof. \\
\end{proof}
The Pro-CoMoM recursive solution is qualitatively similar to the CoMoM recursion. The only significant differences are two: (a) the slightly different structure of (\ref{EQU:PCE})-(\ref{EQU:PPC}) can occasionally result in an over-determined linear system, but this can still be solved, e.g., by the Wiedemann algorithm; (b) for each population vector $\vec N$, Pro-CoMoM needs to evaluate a set of $N+1$ marginal probabilities $P_k(n|\vec N)$ using a different linear system for each of them. This implies that the computational costs of Pro-CoMoM grow as $N+1$ times the computational costs of CoMoM. Since also the MVA algorithm for marginal queue-length probabilities (i.e., the load-dependent MVA) has computational requirements that are $N+1$ times those of the original MVA, it is immediate to conclude that the computational gains of CoMoM over MVA apply also to the Pro-CoMoM algorithm when compared to the load-dependent MVA. For illustration purposes, we show below an example of model that is too computationally expensive to solve with the load-dependent MVA and which can instead be solved exactly and efficiently by the Pro-CoMoM algorithm.

\end{definition}
\begin{figure}[t]
\centering
  % Requires \usepackage{graphicx}
{\includegraphics[width=0.75\columnwidth]{fig6}}
  \caption{{\bf Energy Consumption Case Study}. The application servers $2-3$ can reconfigure dynamically their frequency to minimize energy consumption.}
  \label{FIG:twotier}
\end{figure}

%As remarked in the previous proof, the matrices (\ref{EQU:PCoMoM}) are computed similarly to the CoMoM case except for small differences that occur due to the slightly different structure of (\ref{EQU:PCE})-(\ref{EQU:PPC}). The resulting linear system is typically over-determined, thus the minimal form (\ref{EQU:PCoMoM}) is obtained by pre-multiplying both sides by the transpose of the rectangular coefficient matrix. Because of the latter operation, the sparsity and structure of the involved matrices is not preserved and the fine decomposition discussed for CoMoM does not apply to Pro-CoMoM. This makes the Pro-CoMoM requirements growing as
%$\left[{M+R-1 \choose M}M\right]^3 N^3 \log N$ time and $\left[{M+R-1 \choose M}M\right]^2 N^2 \log N$ space,
%which is a $N$ factor more costly in both requirements than the basic CoMoM implementation. Nevertheless, this is much more efficient than computing marginal queue-length probabilities with existing methods, e.g., with the load-dependent MVA\footnote{We recall that the load-dependent MVA can be used also in the case where all servers have constant processing rate. In this case, the method immediately returns the marginal probabilities $P_k(n\,|\,\vec N)$. However, we show in the final case study that the load-dependent MVA does not enjoy the scalability properties of the Pro-CoMoM algorithm and therefore is limited to the evaluation of small models.} that requires $MRN\prod_{r-1}^R(N_r+1)$ time and space. Comparison with Pro-CoMoM reveals that the computational gains of the technique are qualitatively similar to those of CoMoM with respect to existing normalizing constant algorithms.

%We conclude this section by remarking that, given the probabilistic structure of PCE and PPC, one may consider to use these equations to extend the CoMoM technique to the load-dependent case. However, this could work only for the case of a network including a single load-dependent server $k$, since otherwise the CEs and PCs used in the proof of Theorem \ref{THM:Pmarg} do not apply to the subnetwork without $k$ being correct only for models with constant rate servers \cite{Cas06c}.
%\item the proof of PCEs and PPCs generalizes similarly to the evaluation of class-dependent marginal probabilities $P_k(n_r\,|\,\vec N)$ which describe the frequency of observing $n_r$ jobs in queue $k$. In this case, the PCE becomes
%    $${\mathbf TODO}$$
%    and the PPC becomes
%    $${\mathbf TODO}$$
%\item existing methods such as DAC \cite{DeSL89} and ACAL \cite{LiG93} are able to efficiently compute joint queue-length $\Pr(n_1,n_2,\ldots,n_M\,|\,\vec N)$. We do not address this case because for large populations the number itself of joint probabilities grows as $O(N^M)$ which is approximately the same computational cost of DAC and ACAL. Therefore, the analysis is either prohibitively expensive if the population or number of queues is large or already enjoys nearly-optimal computational algorithms.

\subsection{Energy Consumption Analysis Example}
\label{SEC:perfopt}
We apply the Pro-CoMoM recursive technique to a case study of energy management in Web applications. We evaluate a Web application executing on two heterogeneous application servers and serving a population of requests arriving from a shared communication channel, see Figure~\ref{FIG:twotier}.  The servers are able to increase or decrease their frequency without significant overhead. We assume that the frequency scaling in first approximation does not alter the mean service time perceived by the requests. We also assume that both servers $2-3$ in Figure~\ref{FIG:twotier} implement frequency scaling and we define an energy management policy $f_k(n)$ as a function of the number of jobs $n$ at queue $k$ which specifies the CPU frequency of that server associated to that particular state. In this example, we consider
\begin{equation}
f_k(n)=\begin{cases}
1.00 f^{max}_k & n\geq 15, \\
0.50 f^{max}_k & 5\leq n< 15, \\
0.25 f^{max}_k & n< 5, \\
\end{cases}\nonumber
\end{equation}
where $f^{max}_k$ is the maximum operational frequency of the application server $k$'s CPU.
Table~\ref{TAB:casestudyenergy} gives the service demands for the five service classes considered in the example. The mean think times are $Z_1=10ms$, $Z_2=4ms$, $Z_3=2ms$, $Z_4=5ms$, and $Z_5=3ms$. The total request population is $N=70$ distributed across the classes with a mix $\vec \beta=(0.36,0.21,0.2,0.14,0.09)$, $\beta_r=N_r/N$, $1\leq r\leq 5$. Throughout the example we assume that  on average requests pay a single visit to the servers before completing, a condition that follows by imposing $p_1=p_2=0.25$.
\begin{table}[ht!]
\centering
\caption{Service demands  $D_{k,r}$ [ms] for the model in Figure~\ref{FIG:twotier}}
\label{TAB:casestudyenergy}
\begin{tabular}{|c|cc ccc|}
\hline
\emph{\textrm{Queues} vs {\rm Classes}} & C1 &C2 &C3 &C4 &C5 \\
   \hline
{\rm queue~} 1  &     13  &  35 &   75  &  70 &   55\\
{\rm queue~} 2  &     50  &  59 &   26  &  89 &   15\\
{\rm queue~} 3  &     96  &  23 &   51  &  96 &   16\\
   \hline
\end{tabular}
\end{table}

We begin by showing that the MVA approach cannot evaluate accurately power consumption levels in this example. The savings due to energy management are given by $\Delta P=P_{CPU}(1-f^3)$, where $f$ is the normalized CPU frequency in $[0,1]$, and $P_{CPU}$ is a system-dependent coefficient expressed in Watts \cite{CheDQSWG05}. In order to parameterize the example with realistic values, we set $P_{CPU}=22.7$ Watts as measured in \cite{ElnKR02}.

The marginal probabilities computed by Pro-CoMoM for this example are plotted in Figure~\ref{FIG:twotiermarginals}. Using the energy management policy $f_2(n)$ for server $2$, the savings due to frequency scaling are $\Delta P=15.91$ Watts. Instead, starting from the MVA results, one may assume that the marginal probabilities are uniformly distributed around the mean, with the exception of the state $n=0$ where the probability is by definition one less the utilization of queue $3$. According to this approximation, the energy savings of the application server $2$ predicted by MVA are $\Delta P=11.95$ Watts, with an error of $24.88\%$ with respect to the exact Pro-CoMoM estimate. Using frequency scaling for server $3$, instead, Pro-CoMoM computes a negligible $\Delta P=0.0039$ Watts saving which is $720\%$ smaller than the MVA prediction of a saving of $\Delta P=2.82$ Watts. This shows that both for low utilized and heavily utilized servers, the predictions of MVA are affected by consistent errors that make the estimates unreliable. The Pro-CoMoM algorithm, instead, can provide {\em exact} estimates of the energy savings thus supporting with correct information the energy management decisions.
\begin{figure}[t]
\centering
  % Requires \usepackage{graphicx}
{\includegraphics[width=0.55\columnwidth]{fig7}}
  \caption{{\bf Energy Management Case Study: Marginal Queue-Length Probabilities}. The three curves represent the marginal probabilities $P_k(n)$ computed by Pro-CoMoM for the model in Figure~\ref{FIG:twotier}.}
  \label{FIG:twotiermarginals}
\end{figure}

As an additional experiment, we show the scalability of Pro-CoMoM in evaluating marginal probabilities compared to the load-dependent MVA. For the same experiment discussed above, we were able to solve with CoMoM models where the total population is $N=2\times 70=140$, $N=4\times 70=280$, and $N=7\times 70=490$. In these evaluations, the total memory requirement of Pro-CoMoM was always less than $200$MB, whereas the load-dependent MVA is infeasible in all three cases, having a computational requirement ranging from $20$GB to $40$TB. This confirms that Pro-CoMoM is the only existing algorithm that can be used for the exact evaluation of marginal probabilities.

\begin{figure}[t]
\centering
  % Requires \usepackage{graphicx}
\subfigure[MoM basis $V(\vec N)$ for a model with $M=2$ queues and $R=2$ classes]{\includegraphics[width=0.45\columnwidth]{fig8a}}\quad
\subfigure[CoMoM basis $\Lambda(\vec N)$ for a model with $M=2$ queues and $R=3$ classes]{\includegraphics[width=0.45\columnwidth]{fig8b}}
  \caption{{\bf MoM and CoMoM basis comparison}. The MoM basis $V(\vec N)$ combinatorially increases the number of queue replicas, the CoMoM basis $\Lambda(\vec N)$ combinatorially removes customers from the network.}
  \label{FIG:basecompare}
\end{figure}

\section{Methods of Moments: Cost Comparison}
\label{SEC:duality}
In the previous sections we have focused on the comparison of CoMoM with established solution algorithms. In this section, instead, we compare CoMoM with the linear algebraic computational algorithm recently presented in \cite{Cas06b}, henceforth referred to as the Method of Moments (MoM). In both MoM and CoMoM a basis of normalizing constants is recursively computed using a linear system of CEs and PCs. However, the two methods greatly differ for the choice of the basis of unknowns computed at each recursive step. In the simplest implementation based on an equation similar to (\ref{EQU:comom0}), MoM defines its basis $V(\vec N)$ by combinatorially  increasing the number of queue replicas in the model up to $R$ queues more and considering constants for the populations $\vec N$, $\vec N-\vec 1_1$, $\ldots$, $\vec N-\vec 1_{R-1}$. Instead, CoMoM evaluates normalizing constants on a much larger set of populations (up to $M$ jobs less), but without adding more than one queue replica. Both strategies of MoM and CoMoM lead to an increase of the number of PCs and CEs sufficient to reduce the degrees of freedom $df$ of the analysis to zero.

Figure~\ref{FIG:basecompare} compares the bases of MoM for a model with $M=2$ queues and $R=2$ classes with the basis of CoMoM in a model with $M=2$ and $R=3$.
 Even though the CoMoM model has one class more, the combinatorial structure of the two bases is symmetric. This is a consequence of the fact that, while CoMoM uses conditional higher-order moments of queue-lengths, MoM is based on binomial moments of queue-lengths, and it can be shown that the two sets of moments grow with similar combinatorial structure. We also remark that the basis of CoMoM is quite different from the information carried on by other methods. Focusing on class recursions such as RECAL and RGF, which have a linear recursion on the population that is comparable with that of CoMoM, it is possible to show that RECAL implicitly uses a basis that is similar to that of MoM, but where the number of elements is not fixed with $R$, but grows combinatorially with the current population size $N$. That is, at each recursive step, a new ``level'' is added on top of a basis similar to the one in Figure~\ref{FIG:basecompare}(a). Also RGF has a basis similar to RECAL: RGF performs $R$ recursive steps, one for each service class, in which the basis of normalizing constants is increased by $N_r$ levels after the step for class $r$. These comparison clarifies that the structure of the basis of CoMoM is different with respect to methods in the literature, which instead are more similar to MoM in this respect.

We also note that for each model in Figure~\ref{FIG:basecompare}, MoM evaluates $R$ normalizing constants, while CoMoM $M+1$. Therefore, the two bases $V(\vec N)$ and $\Lambda(\vec N)$ have in general different cardinalities. We prove the following statement.

\begin{figure}[t]
\centering
  % Requires \usepackage{graphicx}
 % \subfigure[$M=3$, $R=2$, order $= 12$, $2$ macro-blocks, $2$ micro-blocks, maximum inverted block size $=9$]{\includegraphics[width=0.2 \textwidth]{btf-32}}\quad\quad
 % \subfigure[$M=3$, $R=3$, order $= 30$, $3$ macro-blocks, $4$ micro-blocks, maximum inverted block size $=9$]{\includegraphics[width=0.2 \textwidth]{btf-33}}\quad\quad
  \subfigure[$M=3$, $R=4$, order $= 60$, $4$ macro-blocks, $8$ micro-blocks, maximum inverted block size $=9$]{\includegraphics[width=0.21 \textwidth]{fig9a}}\quad\quad
  \subfigure[$M=3$, $R=5$, order $= 105$, $4$ macro-blocks, $15$ micro-blocks, maximum inverted block size$=9$]{\includegraphics[width=0.2 \textwidth]{fig9b}}
%  \subfigure[$M=3$, $R=3$]{\includegraphics[width=0.15 \textwidth]{btf-33}}
%  \subfigure[$M=3$, $R=4$]{\includegraphics[width=0.15 \textwidth]{btf-34}}
%  \subfigure[$M=3$, $R=5$]{\includegraphics[width=0.15 \textwidth]{btf-35}}\\
  \caption{{\bf CoMoM Block Triangular Form}. Sparsity structure of the CoMoM coefficient matrix $\comom{\bf A}_{1,1}$ for models with different number of service classes. Points indicate nonzero entries; on the main diagonal, dashed lines indicate macro-blocks, continuous lines identify micro-blocks.}
  \label{FIG:A11CoMoM}
\end{figure}
\begin{theorem}
\label{THM:momvscomom}
{\em
The basis $\Lambda(\vec N)$ of CoMoM has
$
{M+R-1 \choose R}(R-1)
$
normalizing constants less then the basis $V(\vec N)$ of MoM.
}
\end{theorem}
\begin{proof}
MoM evaluates $R$  constants for each models with up to $R$ queue replicas more. The cardinality of its basis $V(\vec N)$ is thus
$
\sum_{l=0}^R {M+l-1\choose l}R={M+R-1 \choose M}(M+R),
$
which by comparison with the size of $V(\vec N)$ given in \cite{Cas06b} proves the theorem.
\end{proof}

The result in (\ref{THM:momvscomom}) proves that the simplest implementation of CoMoM is always much more efficient than the simplest implementation of MoM. However, when considering the BTFs of the two methods, the coefficient matrix $\comom{\bf A}_{1,1}$ of CoMoM has the same identical size of the coefficient matrix ${\bf A}_{1,1}$ of MoM. Table~\ref{TAB:momvscomom} compares the properties of the most efficient MoM and CoMoM implementations based on BTFs.
\begin{table}[h]
\renewcommand{\arraystretch}{1.05}
\centering
\caption{Comparison of MoM and CoMoM}
\begin{tabular}{|l|c|c|}
\hline
{\em Algorithm Property} & {\em CoMoM} & {\em MoM} \\
\hline
~& & \\[-7pt]
Linear system order & $\textstyle \comom{\bf A}_{1,1}: {M+R-1 \choose M}M$ & $\textstyle {\bf A}_{1,1}: {M+R-1 \choose R}R$ \\
\hline
%~& & \\[-5pt]
$H/\comom{H}$ coefficients & $\textstyle \comom{H}: \min\{M,R-1\}$ & $\displaystyle {H}: \min\{M,R\}$ \\
\hline
~& & \\[-5pt]
macro-blocks & $\textstyle \comom{\bf C}_{h}: \comom{H}+1$ & $\textstyle {\bf C}_{h}: {H}$ \\
\hline
Number of micro-blocks & $\textstyle \comom{\bf C}^{(t)}_{h}: {R-1 \choose h}$ & $\textstyle {\bf C}^{(t)}_{h}: {M \choose h}$ \\
\hline
%~& & \\[-5pt]
micro-block order & $\textstyle  \comom{\bf C}^{(t)}_{h}:{M \choose h}M$ & $\textstyle {\bf C}^{(t)}_{h}: {R-1 \choose R-h}R$ \\
\hline
\end{tabular}
\label{TAB:momvscomom}
\end{table}
 The formulas for the micro-blocks order, which are critical for the computational requirements of the two methods, indicate the following rule of thumb for determining which between MoM and CoMoM is more efficient on a model with $M$ queues and $R$ classes:
if $M\geq  R$, then the ${\bf C}^{(t)}_{h}$ micro-blocks are the smallest and MoM should be adopted, otherwise; i.e., for $M\leq R-1$ the $\comom{\bf C}^{(t)}_{h}$ micro-blocks are smaller and CoMoM is the method of choice. These observations are confirmed by figures \ref{FIG:A11CoMoM}-\ref{FIG:A11MoM} which show the fine-grain decomposition of MoM and CoMoM on models with increasing number of service classes. While CoMoM is able to effectively reduce the order of the micro-blocks and minimize the computational requirements, the coefficient matrix of MoM remains composed of a few large diagonal blocks which impose larger costs in the linear system solution. Therefore, MoM does not scale as effectively as CoMoM on models with large number of service classes. This makes CoMoM the best exact algorithm for models with many classes of requests.

\begin{figure}[t]
\centering
  % Requires \usepackage{graphicx}
%  \subfigure[$M=3$, $R=2$, order $= 12$, $2$ macro-blocks, $6$ micro-blocks, maximum inverted block size $=2$]{\includegraphics[width=0.2 \textwidth]{btf-mom-32}}\quad\quad
%  \subfigure[$M=3$, $R=3$, order $= 30$, $3$ macro-blocks, $7$ micro-blocks, maximum inverted block size $=6$]{\includegraphics[width=0.2 \textwidth]{btf-mom-33}}
  \subfigure[$M=3$, $R=4$, order $= 60$, $3$ macro-blocks, $7$ micro-blocks, maximum inverted block size $=12$]{\includegraphics[width=0.2 \textwidth]{fig10a}}\quad\quad
  \subfigure[$M=3$, $R=5$, order $= 105$, $3$ macro-blocks, $7$ micro-blocks, maximum inverted block size $=30$]{\includegraphics[width=0.2 \textwidth]{fig10b}}
  \caption{{\bf MoM Block Triangular Form}. Sparsity structure of $\comom{\bf A}_{1,1}$ in MoM on the models used in Figure~\ref{FIG:A11CoMoM}. As the number of classes $R$ grows, the CoMoM BTF in Figure \ref{FIG:A11CoMoM} generates smaller diagonal blocks that the MoM BTF shown above.}
  \label{FIG:A11MoM}
\end{figure}

% The Figure~suggests that, fixing the number of queues, the basis of MoM and CoMoM have a dual structure with respect to queue addition and customer removal: the queue replica additions in MoM generate a basis $V(\vec N)$ that is structurally similar to the CoMoM basis $\Lambda(\vec N)$ for a model with a class more. More precisely, we find the following theorem.
% \begin{theorem}{\em
%  The bases of MoM and CoMoM have the same combinatorial structure, which implies the same number of normalizing constants, if the following {\em customer-queue duality conditions} apply:
% \begin{equation}
% M_{MoM}=R_{CoMoM},\quad  R_{MoM}-1=M_{CoMoM},
% \end{equation}
% that is, the model evaluated by MoM has $M_{MoM}=R_{CoMoM}-1$ queues and $R_{MoM}=M_{CoMoM}$ classes, being $M_{CoMoM}$ and $R_{CoMoM}$ the number of queues and classes used to define the basis $\Lambda(\vec N)$ of CoMoM.
% }
% \end{theorem}
% \begin{proof}
% The bases of CoMoM is
% $$
% {\Lambda}(\vec N)={\lambda}(\vec 0)\cup \lambda(\vec 1_1)\cup \ldots \cup \lambda(\vec 1_{M_{CoMoM}}),
% $$
% where
% $
% \lambda( \vec v)=\{G( \vec v,\vec N-\vec N)\,|\, \textstyle \sum_{r=1}^{R_{CoMoM}} n_r\leq M_{CoMoM} {\,\rm \wedge\,} n_{R_{CoMoM}}=0\}
% $
% and similarly the basis of MoM is
% $$
% {V}(\vec N)={\nu}(\vec 0)\cup \nu(\vec 1_1)\cup \ldots \cup \nu(\vec 1_{R_{MoM}-1}),
% $$
% with
% $
% \nu( \Delta \vec N)=\{G( \vec v,\vec N-\vec N)\,|\, \textstyle \sum_{k=1}^{M_{MoM}} \Delta m_k\leq R_{MoM}\}.
% $
% In both techniques the recursive step starts using the PCs of class $R$ to compute part of the unknown constants for the new population, e.g., as shown for CoMoM in Figure~\ref{FIG:motivatingCoMoM}(a). Therefore, in both methods the linear system evaluated in the second stage of the recursive step is composed by a replication over models with different number of queue replica (MoM) or different populations (CoMoM) of the linear system
% \begin{align}
% \begin{cases}
% \textstyle G(\vec 1_k,\vec N)-\sum_{r=1}^{R-1} D_{k,r}G(\vec 1_k,\vec N-\vec 1_r)=\textstyle c_k,& 1\leq k\leq M; \nonumber\\
% \textstyle -\sum_{k=1}^M m_k D_{k,r}G(\vec 1_k,\vec N-\vec 1_r)=\textstyle c_r^\prime,& 1\leq r\leq R-1. \nonumber
% \end{cases}
% \end{align}
% where the left-hand side includes all unknown constants of the CEs (\ref{EQU:ce}) and PCs (\ref{EQU:genpc}), while $c_k$ and $c_r^\prime$ are known terms including the values of the normalization constants computed in the initial stage using the PCs of class $R$.
%\end{proof}


\section{Conclusion}
\label{SEC:conclusion}
We have introduced CoMoM, a new efficient algorithm for the solution of multiclass queueing network models widely used in capacity planning and performance evaluation of multiclass systems such as multi-tier architectures. The algorithm solves multiclass performance models several orders of magnitude faster and less memory consuming than MVA and existing methods.
We have also shown that CoMoM generalizes to the recursive computation of marginal queue-length probabilities in multiclass models, which are important to estimate state-dependent performance attributes.

A possible line of future research for CoMoM is the numerical stabilization of the linear recursion (\ref{EQU:comom0}). Preliminary experiments indicate that, if such a result is achieved, it may be possible to solve quickly queueing network models with tens of classes and thousands of requests. This would make exact methods computationally competitive with approximate methods, but without the drawbacks of returning inaccurate solutions. Additionally, it would be interested to study if new approximate methods can be derived from (\ref{EQU:comom0}). %A possible research direction would be to explore if regularization methods \cite{Neu98} could be adapted to the stabilization of~(\ref{EQU:comom0}).
%If such a stabilization cannot be obtained, it would be still interesting to obtain preconditioners for ${\bf A}(\vec N)$ that could make iterative methods such as the Wiedemann algorithm much faster than Gaussian elimination on small models. Finally, it would be interesting to explore the efficiency of the solution of (\ref{EQU:comom0}) with p-adic lifting methods popular in exact algebra \cite{Vil07}.
%{\bf STABILIZATION}


%\begin{table}[ht!]
%\centering
%\caption{Comparison of Time Requirements}
%\label{TAB:comparison}
%\begin{tabular}{ccc|c ccccc}
%\hline
%\hline
% & & & \multicolumn{5}{|c}{ Number of operations} \\
%$M$ & $R$ &   $N$& MVA & Conv & {\sc Recal} & MoM & CoMoM\\
%\hline
%2 & 2 & 1000 & $10^{5}$ & $\bf 10^{4}$ & $10^{12}$ & $10^{8}$ & $10^8$\\
%2 & 4 & 1000 & $10^{9}$ & $\bf 10^{8}$ & $10^{12}$ & $10^{9}$ & $\bf 10^8$\\
%2 & 8 & 1000 & $10^{17}$ & $10^{16}$ & $10^{12}$ & $10^{11}$ & $\bf 10^9$\\
%2 & 10 & 1000 & $10^{21}$ & $10^{20}$ & $10^{12}$ & $10^{12}$ & $\bf 10^9$\\
%\hline
%3 & 2 & 1000 & $10^{6}$ & $\bf 10^{5}$ & $10^{15}$ & $10^{8}$ & $10^8$\\
%3 & 4 & 1000 & $10^{11}$ & $10^{10}$ & $10^{15}$ & $\bf 10^{9}$ & $\bf 10^9$\\
%3 & 8 & 1000 & $10^{18}$ & $10^{17}$ & $10^{15}$ & $10^{12}$ & $\bf 10^{10}$\\
%3 & 10 & 1000 & $10^{22}$ & $10^{21}$ & $10^{15}$ & $10^{13}$ & $\bf 10^{10}$\\
%\hline
%4 & 2 & 1000 & $10^{6}$ & $\bf 10^{5}$ & $10^{17}$ & $10^{8}$ & $10^9$\\
%4 & 4 & 1000 & $10^{11}$ & $\bf 10^{10}$ & $10^{17}$ & $\bf 10^{10}$ & $\bf 10^{10}$\\
%4 & 8 & 1000 & $10^{18}$ & $10^{17}$ & $10^{17}$ & $10^{12}$ & $\bf 10^{11}$\\
%4 & 10 & 1000 & $10^{22}$ & $10^{21}$ & $10^{17}$ & $10^{13}$ & $\bf 10^{11}$\\
%\hline
%\hline
%\end{tabular}
%\end{table}


\bibliographystyle{plain}
%\bibliography{../qntheory,../misc}
\begin{thebibliography}{10}

%\bibitem{BalS96}
%G.~Balbo and G.~Serazzi.
%\newblock Asymptotic analysis of multiclass closed queueing networks: Common  bottlenecks.
%\newblock {\em PEVA}, 26(1):51--72, 1996.

\bibitem{Bar94}
R.~Barrett, M.~Berry, T.~Chan, J.~W. Demmel, J.~Donato, J.~Dongarra, V.~Eijkhour, R.~Pozo, C.~Romine, H.~van der Vorst,
\newblock {\em Templates for the solution of Linear Systems}.
\newblock SIAM, 1994.

\bibitem{BasCMP75}
F.~Baskett, K.~M. Chandy, R.~R. Muntz, and F.~G. Palacios.
\newblock Open, closed, and mixed networks of queues with different classes of customers.
\newblock {\em JACM}, 22(2):248--260, 1975.
%
%\bibitem{BolGMT98}
%G.~Bolch, S.~Greiner, H.~de~Meer, and K.~S. Trivedi.
%\newblock {\em Queueing Networks and Markov Chains}.
%\newblock John Wiley and Sons, 1998.

% \bibitem{Buz73}
% J.~P. Buzen.
% \newblock Computational algorithms for closed queueing networks with
%   exponential servers.
% \newblock {\em Comm. ACM}, 16(9):527--531, 1973.

\bibitem{Cas06b}
G.~Casale.
\newblock An efficient algorithm for the exact analysis of multiclass queueing
  networks with large population sizes.
\newblock In {\em Proc. of ACM SIGMETRICS/IFIP Performance 2006},
  pages 169--180, ACM, 2006.

%\bibitem{Cas06c}
%G.~Casale.
%\newblock On single-class load-dependent normalizing constant equations.
%\newblock In {\em Proc. of QEST 2006}, pages 333--342. IEEE Press, 2006.

\bibitem{Cas07}
G.~Casale.
\newblock The Method of Moments for Queueing Networks.
\newblock (under submission, journal version of \cite{Cas06b}), 2008.

\bibitem{TR}
G.~Casale.
\newblock CoMoM: A Class-Oriented Algorithm for Probabilistic
Evaluation of Multiclass Queueing Networks.
\newblock Tech. Rep. WM-CS-2008-04, Dept. of Computer Science, College of William and Mary, 2008.

%\bibitem{ChaN82}
%K.~M. Chandy and D.~Neuse.
%\newblock Linearizer: {A} heuristic algorithm for queuing network models of
%  computing systems.
%\newblock {\em Comm. ACM}, 25(2):126--134, 1982.

\bibitem{ChaS80}
K.~M. Chandy and C.~H. Sauer.
\newblock Computational algorithms for product-form queueing networks models of
  computing systems.
\newblock {\em Comm. ACM}, 23(10):573--583, 1980.

\bibitem{CheDQSWG05}
Y.~Chen, A.~Das, W.~Qin, A.~Sivasubramaniam, Q.~Wang, and N.~Gautam.
\newblock Managing server energy and operational costs in hosting centers.
\newblock In {\em Proc. of ACM SIGMETRICS 2005}, pages 303--314, 2005.

%\bibitem{ChrGSN08}
%K.~Christensen, C.~Gunaratne, S.~Suen, and B.~Nordman.
%\newblock Reducing the energy consumption of ethernet with adaptive link rate
%  ({ALR}).
%\newblock {\em IEEE Tran. on Computers}, 57(4), 2008.

\bibitem{ChuAT97}
Y.~Chu, C. J.~Antonelli, and T. J.~Teorey.
\newblock Performance Modeling of the PeopleSoft Multi-Tier Remote Computing Architecture.
\newblock University of Michigan, Technical Report CITI 97�5, Dec 1997.

 \bibitem{Coh78}
 D.~J.~A. Cohen.
 \newblock {\em Basic Techniques of Combinatorial Theory}.
 \newblock John Wiley and Sons, 1978.

\bibitem{ConG86}
A.~E. Conway and N.~D. Georganas.
\newblock {\sc RECAL} - {A} new efficient algorithm for the exact analysis of
  multiple-chain closed queueing networks.
\newblock {\em JACM}, 33(4):768--791, 1986.

%\bibitem{DeSL89}
%E.~de~Sousa~e Silva and S.~S. Lavenberg.
%\newblock Calculating joint queue-length distributions in product-form queueing
%  networks.
%\newblock {\em JACM}, 36(1):194--207, 1989.

%\bibitem{DenB78}
%P.~J. Denning and J.~P. Buzen.
%\newblock The operational analysis of queueing network models.
%\newblock {\em ACM Comp. Surv.}, 10(3):225--261, 1978.

\bibitem{ElnKR02}
E.~N. Elnozahy, M.~Kistler, and R.~Rajamony.
\newblock Energy-efficient server clusters.
\newblock In {\em PACS}, pages 179--196, 2002.

\bibitem{HarC02}
P.~G. Harrison and S.~Coury.
\newblock On the asymptotic behaviour of closed multiclass queueing networks.
\newblock {\em PEVA}, 47(2):131--138, 2002.

\bibitem{HarL04}
P.~G. Harrison and T.~T. Lee.
\newblock A new recursive algorithm for computing generating functions in
  closed queueing networks.
\newblock In {\em Proc. of IEEE MASCOTS}, pages
  223--230, IEEE Press, 2004.

\bibitem{KouB03}
S.~Kounev and A.~Buchmann.
\newblock Performance modeling and evaluation of large-scale {J2EE} applications.
\newblock In {\em Proc. of the 29th Int. Conference of the Computer
  Measurement Group {(CMG)}}, 273--283, 2003.

\bibitem{LaMO90}
B. A.~LaMacchia and A. M.~Odlyzko.
\newblock Solving Large Sparse Linear Systems over Finite Fields.
\newblock In {\em Proc. of the 10th Annual International Cryptology Conference on Advances in Cryptology}, 109--133, 1990.

\bibitem{Lam82}
S.~Lam.
\newblock Dynamic scaling and growth behavior of queueing network normalization
  constants.
\newblock {\em JACM}, 29(2):492--513, 1982.

\bibitem{Lam83}
S.~Lam.
\newblock A simple derivation of the {\sc mva} and {\sc lbanc} algorithms from
  the convolution algorithm.
\newblock {\em IEEE T. Comp.}, 32:1062--1064, 1983.

%\bibitem{Lav89}
%S.~S. Lavenberg.
%\newblock A perspective on queueing models of computer performance.
%\newblock {\em PEVA}, 10(1):53--76, 1989.

\bibitem{LazZGS84}
E.~D. Lazowska, J.~Zahorjan, G.~S. Graham, and K.~C. Sevcik.
\newblock {\em Quantitative System Performance}.
\newblock Prentice-Hall, 1984.

%\bibitem{LiG93}
%M.~Li and N.~D. Georganas.
%\newblock {ACAL}-an efficient adaptive chain oriented algorithm for the
%  analysis of multiple chain closed queueing networks.
%\newblock {\em PEVA}, 18(3):225--235, 1993.
%
\bibitem{MenA98}
D.~Menasce and V.~A.~F. Almeida.
\newblock {\em Capacity Planning for Web Performance: Metrics, Models, and
  Methods}.
\newblock Prentice Hall, 1998.

%\bibitem{Neu98}
%A.~Neumaier.
%\newblock {Solving ill-conditioned and singular linear systems:
%		 {A} tutorial on regularization}.
%\newblock {\em SIAM Review}, 40(3):636--666, 1998.

\bibitem{ReiK75b}
M.~Reiser and H.~Kobayashi.
\newblock Queueing networks with multiple closed chains: Theory and
  computational algorithms.
\newblock {\em IBM J. Res. Dev.}, 19(3):283--294, 1975.

\bibitem{ReiL80}
M.~Reiser and S.~S. Lavenberg.
\newblock Mean-value analysis of closed multichain queueing networks.
\newblock {\em JACM}, 27(2):312--322, 1980.

\bibitem{SchS71}
A.~Schonhage and V.~Strassen.
\newblock Schnelle multiplikation großer zahlen.
\newblock {\em Computing}, 7:281--292, Springer Wien, 1971.

%\bibitem{Sch79}
%P.~J. Schweitzer.
%\newblock Approximate analysis of multiclass closed networks of queues.
%\newblock In {\em Int'l Conf. on Stoch. Control and Optim.},  25--29, 1979.

\bibitem{UrgPSST05}
B.~Urgaonkar, G.~Pacifici, P.~J. Shenoy, M.~Spreitzer, and A.~N. Tantawi.
\newblock An analytical model for multi-tier internet services and its
  applications.
\newblock In {\em Proc. of ACM SIGMETRICS 2005}, pages 291--302. ACM, 2005.

\bibitem{Vil07}
G.~Villard.
\newblock Some recent progress in exact linear algebra and related questions
\newblock In {\em Proc. of ISSAC 2007}, pages 391--392. ACM, 2007.

%\bibitem{WanZL08}
%J.~Wang, H.~Zhu, and D.~Li.
%\newblock eraid: Conserving energy in conventional disk-based raid system.
%\newblock {\em IEEE T. Comp.}, 57(4), 2008.

\bibitem{Zah84}
J.~Zahorjan.
\newblock The distribution of network states during residence times in product
  form queueing networks.
\newblock {\em PEVA}, 4:99--104, 1984.

\end{thebibliography}

\section*{\sc Appendix A: Hybrid MVA/CoMoM Algorithm}
In this appendix, we sketch the hybrid MVA/CoMoM algorithm to be used in degenerate models with singular matrix ${\bf \comom{A}}(\vec N)$.
After identifying and removing the set of linearly dependent rows in ${\bf \comom{A}}(\vec N)$ with numerical methods, we are left with an under-determined linear system with $d$ rows less than the original, but we still have to compute all normalizing constants in $\Lambda(\vec N)$ because otherwise we would miss the term ${\bf B}(\vec N)\Lambda(\vec N-\vec 1_R)$ at the next recursive step. We propose to recursive compute these $d$ constants by a hybrid MVA/CoMoM algorithm. In this hybrid algorithm, we replace (\ref{EQU:comom0}) by a more general recursion
\[
\textstyle {\bf \comom{A}^\ast}(\vec N){\Lambda^\ast}(\vec N)=\sum_{r\in {\cal R}^\ast} \comom{{\bf B}_r^\ast}(\vec N){\Lambda^\ast}(\vec N-1_r)+\comom{{\bf B}_R^\ast}(\vec N){\Lambda^\ast}(\vec N-1_R),
\]
where ${\Lambda^\ast}(\vec N)\subset {\Lambda}(\vec N)$ does not longer include the $d$ normalizing constants, $\det({\bf \comom{A}^\ast}(\vec N))\neq 0$, and ${\cal R}^\ast$ is the set of classes on which we perform the additional recursions. Similarly to (\ref{EQU:comom0}), the choice of the normalizing constants in ${\Lambda^\ast}(\vec N)$ immediately defines from (\ref{EQU:ce})-(\ref{EQU:pc}) the structure of the matrices ${\bf \comom{A}^\ast}(\vec N)$, $\comom{{\bf B}_r^\ast}(\vec N)$, and $\comom{{\bf B}_R^\ast}(\vec N)$. For example, suppose that in a degenerate model with $R=4$ classes we remove the normalizing constant $G(\vec N-1_2)$ from ${\Lambda}(\vec N)$. Then, the recursion on class $4$~would be
\begin{equation}
\textstyle {\bf \comom{A}^\ast}(\vec N){\Lambda^\ast}(\vec N)=\comom{{\bf B}^\ast_2}(\vec N){\Lambda^\ast}(\vec N-1_2)+\comom{{\bf B}^\ast_4}(\vec N){\Lambda^\ast}(\vec N-1_4),\nonumber
\end{equation}
where we are now obtaining the value of $G(\vec N-1_2)$ from ${\Lambda^\ast}(\vec N-1_2)$ instead than from $\Lambda^\ast(\vec N)$. It is also possible to note that ${\Lambda^\ast}(\vec N-1_2)$ includes other elements of ${\Lambda^\ast}(\vec N)$, e.g., $G(\vec 1_k, \vec N-1_2)$, $1\leq k \leq M$. This redundancy can be exploited to minimize computational costs of the hybrid method: we first select the $d$ normalizing constants to remove from ${\Lambda^\ast}(\vec N)$ as a set of unknowns such that $\det({\bf \comom{A}^\ast}(\vec N))\neq 0$ and which minimizes the number of elements in ${\cal R}^\ast$. Then, we remove all redundant constants by defining ${\Lambda^\ast}(\vec N)$ as the basis of a model with classes $\{1,2,\ldots,R\}/{\cal R}^\ast$. Note that ${\bf \comom{A}^\ast}(\vec N)$ becomes ${\bf \comom{A}}(\vec N)$ for a model with $R-card({\cal R}^\ast)$ classes and thus we can still apply the BTFs defined in Section~\ref{SEC:btf}. The computational cost of the hybrid algorithm can be computed in advance as the cost of solving with CoMoM a model with $R-card({\cal R}^\ast)$ classes multiplied by the number of the additional populations spanned by the MVA-like recursion, which is $\prod_{r\in {\cal R}^\ast} (N_r+1)$. %This allows to estimate in advance the computational costs of the hybrid algorithm.

\begin{biography}[{\includegraphics[width=1in,height=1.25in,clip,keepaspectratio]{photo.eps}}]{Giuliano Casale} received the MS and the PhD degrees in computer engineering from the Politecnico di Milano, Milan, Italy, in 2002 and 2006, respectively. From January 2007 he is a postdoctoral research associate at the College of William and Mary, Williamsburg, Virginia, where he studies the performance impact of burstiness in systems. In Fall 2004, he was a visiting scholar at UCLA studying bounds for queueing networks. His research interests include performance evaluation, modeling, optimization, and simulation. He is member of the ACM, IEEE, and officer of the IEEE Hampton Roads Section.
\end{biography}

%\appendix
%\noindent {\sc CoMoM Matrices Computation}
\end{document}

%% TODO CAMERA READY
%DONE - REPLACE QUMOM -> MOM
%DONE - dire che il modello di peoplesoft e' SIMILE a quello del TR, manca il LD
%DONE - remove last paragraph of DVS example with the "few hours`` thing
%DONE - rimuovi Estimates refer to an architecture that can process $10^8$ operations per second.
%DONE - RGF cannot solve model with Z ->  change Kounev model to Zr=0 for all classes
%DONE - add ceil to n=N\log(Dmax(M+R+N))
%DONE - RGF with Z?
%DONE - rename base to basis
%DONE - remove DVS?
%DONE - COMOM requirements: remove the value of N=5000?
%NOTDONE - usa una notazione diverse da lambda per la definizione di base perche' nel motivating example lambda e' un'altra roba
% - RECHECK l'interpretazione probabilistica per G(N-c_k) sembra sbagliata, come fa a dipendere da k?
